% ----------------------------------------------------------------------
%  1.2.1. Кристаллографические проекции
% ----------------------------------------------------------------------
\begin{orientir}
Раздел даёт рабочий аппарат кристаллографа: как изобразить на плоскости
ориентацию граней и направлений (проекции, 1.2.1), как описать решётку
(1.2.2), как обозначать узлы, направления и плоскости (индексы Миллера,
1.2.3), что такое обратная решётка (1.2.4), как переходить между
системами координат (1.2.5--1.2.6) и как считать углы и расстояния
(1.2.7). Завершает раздел ячейка Вигнера--Зейтца (1.2.8).
\end{orientir}

\subsection{Кристаллографические проекции}\label{ssec:projections}

Для многих задач важна не форма кристалла, а только \emph{взаимная
ориентация} его граней и рёбер. Поэтому вместо кристалла рассматривают
\term{кристаллографический комплекс}: все параллельные грани заменяют
одной плоскостью, все параллельные рёбра~— одной прямой, и все они
проходят через одну точку~— \emph{центр комплекса}. Так, кубооктаэдр (куб
со срезанными вершинами) превращается в три взаимно перпендикулярные
плоскости граней куба и четыре наклонные плоскости срезов.

\begin{opredelenie}[Зона]
Совокупность плоскостей, пересекающихся по общей прямой, называют
\term{кристаллографической зоной}, а эту прямую~— \term{осью зоны}.
Комплекс состоит из ряда зон.
\end{opredelenie}

Часто удобнее \term{обратный} (полярный) комплекс: каждую плоскость
заменяют нормалью к ней, а каждую прямую~— перпендикулярной ей плоскостью.
Зона в обратном комплексе~— это пучок нормалей, лежащих в одной
плоскости, перпендикулярной оси зоны (рис.~\ref{fig:zone}). Для
кубического кристалла прямой и обратный комплексы совпадают.

\begin{figure}[ht]
\centering
\begin{tikzpicture}[x={(1cm,0cm)}, y={(0.45cm,0.32cm)}, z={(0cm,1cm)}, scale=1.05]
  \draw[fmRed, very thick] (0,0,-1.9) -- (0,0,2.1) node[above, font=\small]{ось зоны};
  \foreach \a/\c in {0/fmBlue, 60/fmGreen, 120/fmAmber}{
    \pgfmathsetmacro{\cx}{1.6*cos(\a)}\pgfmathsetmacro{\cy}{1.6*sin(\a)}
    \fill[\c, opacity=0.18] (\cx,\cy,-1.5) -- (\cx,\cy,1.5) -- (-\cx,-\cy,1.5) -- (-\cx,-\cy,-1.5) -- cycle;
    \draw[\c, thin] (\cx,\cy,-1.5) -- (\cx,\cy,1.5) -- (-\cx,-\cy,1.5) -- (-\cx,-\cy,-1.5) -- cycle;
    \pgfmathsetmacro{\nx}{1.5*cos(\a+90)}\pgfmathsetmacro{\ny}{1.5*sin(\a+90)}
    \draw[\c, thick, -{Stealth[length=2mm]}] (0,0,0) -- (\nx,\ny,0);}
  \draw[fmGray, dashed] (0,0,0) circle[radius=1.5];
  \fill (0,0,0) circle (1.3pt);
  \node[font=\scriptsize, fmGray, anchor=west] at (1.5,-1.6,0) {плоскость нормалей};
\end{tikzpicture}
\caption{Кристаллографическая зона: три плоскости, пересекающиеся по оси
зоны (красная линия), и их нормали (стрелки). Нормали всех плоскостей
зоны лежат в одной плоскости, перпендикулярной оси зоны,~— так зона
выглядит в обратном комплексе}
\label{fig:zone}
\end{figure}

\paragraph{Виды проекций}
Комплекс нужно изобразить на плоскости. Способы различаются тем, какой
комплекс проектируют (прямой или обратный) и на что
(табл.~\ref{tab:projections}).

\begin{table}[ht]
\centering
\caption{Кристаллографические проекции}\label{tab:projections}
\small
\begin{tabularx}{\textwidth}{@{}>{\raggedright\arraybackslash\hspace{0pt}}p{33mm}>{\raggedright\arraybackslash\hspace{0pt}}p{25mm}>{\raggedright\arraybackslash\hspace{0pt}}X>{\raggedright\arraybackslash\hspace{0pt}}X>{\raggedright\arraybackslash\hspace{0pt}}X@{}}
\toprule
Проекция & Что и куда & Плоскость & Направление & Зона \\
\midrule
линейная & прямой $\to$ плоскость & прямая & точка & прямые через точку \\
гномоническая & обратный $\to$ плоскость & точка & прямая & точки на прямой \\
сферическая & прямой $\to$ сфера & большой круг & две точки & большие круги через две точки \\
гномосферическая & обратный $\to$ сфера & две точки & большой круг & точки на большом круге \\
\textbf{стереографическая} & прямой $\to$ плоскость & дуга большого круга & точка & дуги через две точки \\
\textbf{гномостереографическая} & обратный $\to$ плоскость & точка (полюс) & дуга большого круга & точки на одной дуге \\
\bottomrule
\end{tabularx}
\end{table}

Линейной и гномонической проекциям для изображения всего комплекса нужна
бесконечная плоскость, а сферические проекции неплоские. На практике
используют почти исключительно \emph{стереографическую} и
\emph{гномостереографическую} проекции. Гномоническую иногда применяют
при обработке лауэграмм.

\paragraph{Стереографическая проекция}
Центр комплекса помещают в центр $O$ сферы; плоскость проекции проходит
через $O$ (рис.~\ref{fig:stereo}\,\textit{а}). Её сечение сферой~—
\term{основной круг} проекции, концы перпендикулярного ей диаметра~—
\emph{полюсы} $N$ (над плоскостью) и $S$ (под ней). Прямую $L$ комплекса
продолжают до пересечения со сферой в точке $M_1$ и соединяют $M_1$ с
южным полюсом $S$; точка $M_1'$, где этот луч пересекает плоскость
проекции,~— \emph{стереографическая проекция прямой}. Второй выход той же
прямой, $M_2$ в южном полушарии, проектируют из $N$ и отмечают
\emph{крестиком} $M_2'$.

\begin{figure}[ht]
\centering
\begin{subfigure}[b]{0.52\textwidth}\centering
\begin{tikzpicture}[scale=1.0]
  \def\R{2.2}
  \draw[atomLine, thick] (0,0) circle (\R);
  \draw[fmGray, thick] (-2.9,0) -- (2.9,0) node[right, font=\small]{$P$};
  \draw[fmGray, dashed] (0,-\R) -- (0,\R);
  \fill (0,\R) circle (1.4pt) node[above, font=\small]{$N$};
  \fill (0,-\R) circle (1.4pt) node[below, font=\small]{$S$};
  \fill (0,0) circle (1.2pt) node[below left, font=\small]{$O$};
  \coordinate (M1) at (50:\R); \coordinate (M2) at (230:\R);
  \draw[fmBlue, very thick] (M2) -- (M1) node[pos=0.88, right, font=\small]{$L$};
  \fill[fmBlue] (M1) circle (1.6pt) node[above right, font=\small]{$M_1$};
  \fill[fmBlue] (M2) circle (1.6pt) node[below left, font=\small]{$M_2$};
  \draw[fmRed] (0,-\R) -- (M1);
  \draw[fmRed, dashed] (0,\R) -- (M2);
  \fill[fmRed] (0.8007,0) circle (1.8pt) node[below right=-1pt, font=\small]{$M_1'$};
  \node[fmRed, font=\small] at (-0.8007,0) {$\times$};
  \node[fmRed, font=\small, above left=-1pt] at (-0.8007,0) {$M_2'$};
\end{tikzpicture}
\caption{}
\end{subfigure}\hfill
\begin{subfigure}[b]{0.44\textwidth}\centering
\begin{tikzpicture}[scale=1.0]
  \def\R{2.2}
  \draw[atomLine, thick] (0,0) circle (\R);
  \fill (0,0) circle (1.2pt);
  \fill[fmRed] (0.8007,0) circle (1.8pt) node[above right=-1pt, font=\small]{$M_1'$};
  \node[fmRed, font=\small] at (-0.8007,0) {$\times$};
  \node[fmRed, font=\small, above left=-1pt] at (-0.8007,0) {$M_2'$};
  \begin{scope}
    \clip (0,0) circle (\R);
    \draw[fmBlue, thick] (-2.622,0) circle (3.4227);
    \draw[fmBlue, thick, dashed] (2.622,0) circle (3.4227);
  \end{scope}
  \node[fmBlue, font=\small] at (1.05,1.25) {$p$};
  \node[fmBlue, font=\small] at (-1.0,1.3) {$p'$};
\end{tikzpicture}
\caption{}
\end{subfigure}
\caption{Стереографическая проекция: \textit{а}~— построение (сечение
сферы плоскостью, содержащей ось $NS$): выход прямой $L$ в северном
полушарии $M_1$ проектируется из $S$ в точку $M_1'$, выход в южном
полушарии $M_2$~— из $N$ в крестик $M_2'$; \textit{б}~— вид на плоскость
проекции: прямая изображается точкой и крестиком на равных расстояниях от
центра, наклонная плоскость~— дугой большого круга $p$ (её часть в южном
полушарии~— штриховой дугой $p'$)}
\label{fig:stereo}
\end{figure}

\begin{keyeq}\label{eq:stereo-rho}
\rho=r\,\operatorname{tg}\frac\theta2
\end{keyeq}
где $\theta$~— угол между прямой и осью $ON$, $r$~— радиус основного круга,
$\rho$~— расстояние проекции от центра. Прямая, перпендикулярная
плоскости проекции ($\theta=0$), проектируется в центр; прямая, лежащая в
плоскости проекции ($\theta=90^\circ$),~— в две точки основного круга.

Плоскость комплекса проектируется дугой большого круга, опирающейся на
диаметр основного круга (рис.~\ref{fig:stereo}\,\textit{б}): сплошной~— её
часть над плоскостью проекции, штриховой~— под ней. Вертикальная
плоскость изображается диаметром, плоскость проекции~— основным кругом.
Зона~— серия дуг, проходящих через две диаметрально противоположные
точки~— проекции оси зоны.

\begin{glavnoe}[Два свойства, ради которых всё это делается]
\begin{enumerate}
\item Любой круг на сфере (большой или малый) изображается на проекции
окружностью или её дугой.
\item Проекция \emph{конформна}: углы между кривыми на сфере равны углам
между их изображениями. Стереографическое проектирование из $S$
устанавливает взаимно однозначное соответствие между точками сферы (кроме
самой $S$) и плоскости.
\end{enumerate}
\end{glavnoe}

\paragraph{Гномостереографическая проекция}
Это стереографическая проекция обратного комплекса: плоскость заменяют
нормалью $\vb n$ и проектируют точку выхода нормали на сферу. Плоскость
изображается одной точкой~— \term{полюсом}, направление~— дугой большого
круга, а все плоскости одной зоны~— полюсами, лежащими на одной дуге. Эта
проекция удобнее для практических задач, и дальше полюсы плоскостей и
выходы направлений мы будем наносить на один чертёж.

\paragraph{Сетки Вульфа и Болдырева}
Измерять углы на проекции позволяют координатные сетки
(рис.~\ref{fig:nets}). \term{Сетка Вульфа}~— стереографическая проекция
двух семейств кругов на сфере:
\begin{itemize}
\item \emph{меридианы}~— большие круги, проходящие через ось $y$, лежащую в
плоскости проекции, повёрнутые друг относительно друга через $2^\circ$
(на рисунке~— через $10^\circ$). Каждый меридиан~— проекция плоскости, а
все вместе~— проекция зоны с осью $y$;
\item \emph{параллели}~— малые круги, проекции соосных конусов с осью $y$,
полураствор которых меняется через $2^\circ$.
\end{itemize}
Угловая шкала по основному кругу равномерна, по экватору и нулевому
меридиану~— неравномерна, но одинакова. \term{Сетка Болдырева}~—
концентрические окружности (проекции соосных конусов с осью $z$,
перпендикулярной плоскости проекции; радиусы $r\operatorname{tg}(\theta/2)$) и
радиальные прямые (проекции плоскостей зоны с осью $z$).

\begin{figure}[ht]
\centering
\tikzset{wulffA/.style={fmBlue!45, line width=0.25pt},
  wulffB/.style={fmBlue!80, line width=0.55pt},
  wulffC/.style={atomLine, line width=0.9pt}}
\begin{subfigure}[b]{0.48\textwidth}\centering
\begin{tikzpicture}[scale=2.9]
  % создано tools/gen_stereo.py: сетка Вульфа (r = 1)
\draw[wulffA] (0.0000,-1.0000) -- (0.0342,-0.9934) -- (0.0679,-0.9856) -- (0.1011,-0.9768) -- (0.1338,-0.9669) -- (0.1660,-0.9560) -- (0.1976,-0.9441) -- (0.2286,-0.9312) -- (0.2591,-0.9174) -- (0.2888,-0.9026) -- (0.3179,-0.8870) -- (0.3464,-0.8706) -- (0.3741,-0.8533) -- (0.4012,-0.8352) -- (0.4275,-0.8164) -- (0.4531,-0.7968) -- (0.4779,-0.7766) -- (0.5020,-0.7557) -- (0.5252,-0.7341) -- (0.5477,-0.7119) -- (0.5695,-0.6891) -- (0.5904,-0.6658) -- (0.6105,-0.6419) -- (0.6297,-0.6175) -- (0.6482,-0.5927) -- (0.6658,-0.5673) -- (0.6826,-0.5416) -- (0.6986,-0.5154) -- (0.7137,-0.4888) -- (0.7280,-0.4619) -- (0.7414,-0.4346) -- (0.7539,-0.4071) -- (0.7656,-0.3792) -- (0.7765,-0.3510) -- (0.7865,-0.3227) -- (0.7956,-0.2940) -- (0.8039,-0.2652) -- (0.8112,-0.2362) -- (0.8178,-0.2070) -- (0.8234,-0.1777) -- (0.8282,-0.1483) -- (0.8321,-0.1188) -- (0.8352,-0.0891) -- (0.8374,-0.0595) -- (0.8387,-0.0297) -- (0.8391,0.0000) -- (0.8387,0.0297) -- (0.8374,0.0595) -- (0.8352,0.0891) -- (0.8321,0.1188) -- (0.8282,0.1483) -- (0.8234,0.1777) -- (0.8178,0.2070) -- (0.8112,0.2362) -- (0.8039,0.2652) -- (0.7956,0.2940) -- (0.7865,0.3227) -- (0.7765,0.3510) -- (0.7656,0.3792) -- (0.7539,0.4071) -- (0.7414,0.4346) -- (0.7280,0.4619) -- (0.7137,0.4888) -- (0.6986,0.5154) -- (0.6826,0.5416) -- (0.6658,0.5673) -- (0.6482,0.5927) -- (0.6297,0.6175) -- (0.6105,0.6419) -- (0.5904,0.6658) -- (0.5695,0.6891) -- (0.5477,0.7119) -- (0.5252,0.7341) -- (0.5020,0.7557) -- (0.4779,0.7766) -- (0.4531,0.7968) -- (0.4275,0.8164) -- (0.4012,0.8352) -- (0.3741,0.8533) -- (0.3464,0.8706) -- (0.3179,0.8870) -- (0.2888,0.9026) -- (0.2591,0.9174) -- (0.2286,0.9312) -- (0.1976,0.9441) -- (0.1660,0.9560) -- (0.1338,0.9669) -- (0.1011,0.9768) -- (0.0679,0.9856) -- (0.0342,0.9934) -- (0.0000,1.0000);
\draw[wulffA] (0.0000,-1.0000) -- (0.0324,-0.9876) -- (0.0640,-0.9743) -- (0.0948,-0.9602) -- (0.1248,-0.9453) -- (0.1540,-0.9296) -- (0.1824,-0.9132) -- (0.2100,-0.8961) -- (0.2367,-0.8784) -- (0.2626,-0.8601) -- (0.2877,-0.8413) -- (0.3120,-0.8219) -- (0.3355,-0.8020) -- (0.3582,-0.7816) -- (0.3801,-0.7608) -- (0.4012,-0.7396) -- (0.4216,-0.7179) -- (0.4411,-0.6959) -- (0.4599,-0.6736) -- (0.4779,-0.6509) -- (0.4952,-0.6280) -- (0.5117,-0.6047) -- (0.5274,-0.5812) -- (0.5425,-0.5575) -- (0.5568,-0.5335) -- (0.5704,-0.5093) -- (0.5833,-0.4850) -- (0.5955,-0.4604) -- (0.6069,-0.4357) -- (0.6177,-0.4108) -- (0.6278,-0.3857) -- (0.6373,-0.3606) -- (0.6460,-0.3353) -- (0.6541,-0.3099) -- (0.6615,-0.2844) -- (0.6683,-0.2588) -- (0.6743,-0.2332) -- (0.6798,-0.2074) -- (0.6846,-0.1816) -- (0.6887,-0.1558) -- (0.6923,-0.1299) -- (0.6951,-0.1040) -- (0.6973,-0.0780) -- (0.6989,-0.0520) -- (0.6999,-0.0260) -- (0.7002,0.0000) -- (0.6999,0.0260) -- (0.6989,0.0520) -- (0.6973,0.0780) -- (0.6951,0.1040) -- (0.6923,0.1299) -- (0.6887,0.1558) -- (0.6846,0.1816) -- (0.6798,0.2074) -- (0.6743,0.2332) -- (0.6683,0.2588) -- (0.6615,0.2844) -- (0.6541,0.3099) -- (0.6460,0.3353) -- (0.6373,0.3606) -- (0.6278,0.3857) -- (0.6177,0.4108) -- (0.6069,0.4357) -- (0.5955,0.4604) -- (0.5833,0.4850) -- (0.5704,0.5093) -- (0.5568,0.5335) -- (0.5425,0.5575) -- (0.5274,0.5812) -- (0.5117,0.6047) -- (0.4952,0.6280) -- (0.4779,0.6509) -- (0.4599,0.6736) -- (0.4411,0.6959) -- (0.4216,0.7179) -- (0.4012,0.7396) -- (0.3801,0.7608) -- (0.3582,0.7816) -- (0.3355,0.8020) -- (0.3120,0.8219) -- (0.2877,0.8413) -- (0.2626,0.8601) -- (0.2367,0.8784) -- (0.2100,0.8961) -- (0.1824,0.9132) -- (0.1540,0.9296) -- (0.1248,0.9453) -- (0.0948,0.9602) -- (0.0640,0.9743) -- (0.0324,0.9876) -- (0.0000,1.0000);
\draw[wulffB] (0.0000,-1.0000) -- (0.0297,-0.9823) -- (0.0584,-0.9639) -- (0.0860,-0.9451) -- (0.1127,-0.9258) -- (0.1384,-0.9061) -- (0.1631,-0.8860) -- (0.1869,-0.8656) -- (0.2098,-0.8448) -- (0.2318,-0.8238) -- (0.2529,-0.8025) -- (0.2732,-0.7809) -- (0.2927,-0.7592) -- (0.3114,-0.7372) -- (0.3293,-0.7151) -- (0.3464,-0.6928) -- (0.3628,-0.6704) -- (0.3785,-0.6479) -- (0.3934,-0.6253) -- (0.4077,-0.6025) -- (0.4213,-0.5797) -- (0.4342,-0.5568) -- (0.4465,-0.5339) -- (0.4582,-0.5109) -- (0.4692,-0.4879) -- (0.4797,-0.4648) -- (0.4896,-0.4416) -- (0.4988,-0.4185) -- (0.5076,-0.3953) -- (0.5157,-0.3721) -- (0.5234,-0.3489) -- (0.5305,-0.3257) -- (0.5370,-0.3025) -- (0.5431,-0.2792) -- (0.5486,-0.2560) -- (0.5537,-0.2327) -- (0.5582,-0.2094) -- (0.5622,-0.1862) -- (0.5658,-0.1629) -- (0.5689,-0.1396) -- (0.5715,-0.1164) -- (0.5736,-0.0931) -- (0.5752,-0.0698) -- (0.5764,-0.0465) -- (0.5771,-0.0233) -- (0.5774,0.0000) -- (0.5771,0.0233) -- (0.5764,0.0465) -- (0.5752,0.0698) -- (0.5736,0.0931) -- (0.5715,0.1164) -- (0.5689,0.1396) -- (0.5658,0.1629) -- (0.5622,0.1862) -- (0.5582,0.2094) -- (0.5537,0.2327) -- (0.5486,0.2560) -- (0.5431,0.2792) -- (0.5370,0.3025) -- (0.5305,0.3257) -- (0.5234,0.3489) -- (0.5157,0.3721) -- (0.5076,0.3953) -- (0.4988,0.4185) -- (0.4896,0.4416) -- (0.4797,0.4648) -- (0.4692,0.4879) -- (0.4582,0.5109) -- (0.4465,0.5339) -- (0.4342,0.5568) -- (0.4213,0.5797) -- (0.4077,0.6025) -- (0.3934,0.6253) -- (0.3785,0.6479) -- (0.3628,0.6704) -- (0.3464,0.6928) -- (0.3293,0.7151) -- (0.3114,0.7372) -- (0.2927,0.7592) -- (0.2732,0.7809) -- (0.2529,0.8025) -- (0.2318,0.8238) -- (0.2098,0.8448) -- (0.1869,0.8656) -- (0.1631,0.8860) -- (0.1384,0.9061) -- (0.1127,0.9258) -- (0.0860,0.9451) -- (0.0584,0.9639) -- (0.0297,0.9823) -- (0.0000,1.0000);
\draw[wulffA] (0.0000,-1.0000) -- (0.0261,-0.9775) -- (0.0511,-0.9548) -- (0.0750,-0.9319) -- (0.0979,-0.9090) -- (0.1197,-0.8859) -- (0.1405,-0.8628) -- (0.1604,-0.8397) -- (0.1794,-0.8166) -- (0.1975,-0.7935) -- (0.2148,-0.7703) -- (0.2313,-0.7473) -- (0.2470,-0.7242) -- (0.2620,-0.7012) -- (0.2763,-0.6783) -- (0.2899,-0.6554) -- (0.3028,-0.6326) -- (0.3151,-0.6098) -- (0.3268,-0.5872) -- (0.3379,-0.5646) -- (0.3484,-0.5421) -- (0.3584,-0.5196) -- (0.3679,-0.4973) -- (0.3768,-0.4750) -- (0.3853,-0.4528) -- (0.3932,-0.4307) -- (0.4007,-0.4087) -- (0.4077,-0.3867) -- (0.4143,-0.3648) -- (0.4204,-0.3430) -- (0.4262,-0.3212) -- (0.4315,-0.2995) -- (0.4364,-0.2778) -- (0.4409,-0.2563) -- (0.4450,-0.2347) -- (0.4488,-0.2132) -- (0.4521,-0.1918) -- (0.4551,-0.1704) -- (0.4578,-0.1490) -- (0.4601,-0.1277) -- (0.4620,-0.1063) -- (0.4635,-0.0850) -- (0.4647,-0.0638) -- (0.4656,-0.0425) -- (0.4661,-0.0212) -- (0.4663,0.0000) -- (0.4661,0.0212) -- (0.4656,0.0425) -- (0.4647,0.0638) -- (0.4635,0.0850) -- (0.4620,0.1063) -- (0.4601,0.1277) -- (0.4578,0.1490) -- (0.4551,0.1704) -- (0.4521,0.1918) -- (0.4488,0.2132) -- (0.4450,0.2347) -- (0.4409,0.2563) -- (0.4364,0.2778) -- (0.4315,0.2995) -- (0.4262,0.3212) -- (0.4204,0.3430) -- (0.4143,0.3648) -- (0.4077,0.3867) -- (0.4007,0.4087) -- (0.3932,0.4307) -- (0.3853,0.4528) -- (0.3768,0.4750) -- (0.3679,0.4973) -- (0.3584,0.5196) -- (0.3484,0.5421) -- (0.3379,0.5646) -- (0.3268,0.5872) -- (0.3151,0.6098) -- (0.3028,0.6326) -- (0.2899,0.6554) -- (0.2763,0.6783) -- (0.2620,0.7012) -- (0.2470,0.7242) -- (0.2313,0.7473) -- (0.2148,0.7703) -- (0.1975,0.7935) -- (0.1794,0.8166) -- (0.1604,0.8397) -- (0.1405,0.8628) -- (0.1197,0.8859) -- (0.0979,0.9090) -- (0.0750,0.9319) -- (0.0511,0.9548) -- (0.0261,0.9775) -- (0.0000,1.0000);
\draw[wulffA] (0.0000,-1.0000) -- (0.0218,-0.9734) -- (0.0426,-0.9470) -- (0.0622,-0.9208) -- (0.0808,-0.8949) -- (0.0985,-0.8692) -- (0.1153,-0.8438) -- (0.1312,-0.8186) -- (0.1463,-0.7937) -- (0.1606,-0.7690) -- (0.1742,-0.7446) -- (0.1871,-0.7204) -- (0.1993,-0.6965) -- (0.2109,-0.6728) -- (0.2219,-0.6494) -- (0.2324,-0.6262) -- (0.2423,-0.6032) -- (0.2516,-0.5804) -- (0.2605,-0.5578) -- (0.2689,-0.5355) -- (0.2769,-0.5133) -- (0.2844,-0.4913) -- (0.2914,-0.4695) -- (0.2981,-0.4479) -- (0.3044,-0.4264) -- (0.3103,-0.4051) -- (0.3159,-0.3839) -- (0.3211,-0.3629) -- (0.3259,-0.3420) -- (0.3304,-0.3212) -- (0.3347,-0.3006) -- (0.3386,-0.2801) -- (0.3422,-0.2596) -- (0.3455,-0.2393) -- (0.3485,-0.2190) -- (0.3512,-0.1989) -- (0.3537,-0.1788) -- (0.3559,-0.1587) -- (0.3578,-0.1388) -- (0.3594,-0.1189) -- (0.3608,-0.0990) -- (0.3620,-0.0791) -- (0.3628,-0.0593) -- (0.3635,-0.0395) -- (0.3638,-0.0198) -- (0.3640,0.0000) -- (0.3638,0.0198) -- (0.3635,0.0395) -- (0.3628,0.0593) -- (0.3620,0.0791) -- (0.3608,0.0990) -- (0.3594,0.1189) -- (0.3578,0.1388) -- (0.3559,0.1587) -- (0.3537,0.1788) -- (0.3512,0.1989) -- (0.3485,0.2190) -- (0.3455,0.2393) -- (0.3422,0.2596) -- (0.3386,0.2801) -- (0.3347,0.3006) -- (0.3304,0.3212) -- (0.3259,0.3420) -- (0.3211,0.3629) -- (0.3159,0.3839) -- (0.3103,0.4051) -- (0.3044,0.4264) -- (0.2981,0.4479) -- (0.2914,0.4695) -- (0.2844,0.4913) -- (0.2769,0.5133) -- (0.2689,0.5355) -- (0.2605,0.5578) -- (0.2516,0.5804) -- (0.2423,0.6032) -- (0.2324,0.6262) -- (0.2219,0.6494) -- (0.2109,0.6728) -- (0.1993,0.6965) -- (0.1871,0.7204) -- (0.1742,0.7446) -- (0.1606,0.7690) -- (0.1463,0.7937) -- (0.1312,0.8186) -- (0.1153,0.8438) -- (0.0985,0.8692) -- (0.0808,0.8949) -- (0.0622,0.9208) -- (0.0426,0.9470) -- (0.0218,0.9734) -- (0.0000,1.0000);
\draw[wulffB] (0.0000,-1.0000) -- (0.0169,-0.9701) -- (0.0329,-0.9407) -- (0.0479,-0.9120) -- (0.0621,-0.8838) -- (0.0755,-0.8561) -- (0.0881,-0.8289) -- (0.1000,-0.8022) -- (0.1113,-0.7760) -- (0.1219,-0.7503) -- (0.1319,-0.7250) -- (0.1414,-0.7001) -- (0.1504,-0.6756) -- (0.1589,-0.6515) -- (0.1669,-0.6277) -- (0.1745,-0.6043) -- (0.1816,-0.5813) -- (0.1884,-0.5585) -- (0.1948,-0.5361) -- (0.2008,-0.5140) -- (0.2065,-0.4921) -- (0.2118,-0.4705) -- (0.2169,-0.4491) -- (0.2216,-0.4280) -- (0.2261,-0.4071) -- (0.2303,-0.3864) -- (0.2342,-0.3659) -- (0.2379,-0.3456) -- (0.2413,-0.3255) -- (0.2445,-0.3055) -- (0.2474,-0.2857) -- (0.2502,-0.2660) -- (0.2527,-0.2465) -- (0.2550,-0.2271) -- (0.2571,-0.2078) -- (0.2590,-0.1886) -- (0.2608,-0.1695) -- (0.2623,-0.1504) -- (0.2636,-0.1315) -- (0.2648,-0.1126) -- (0.2658,-0.0937) -- (0.2665,-0.0749) -- (0.2672,-0.0562) -- (0.2676,-0.0374) -- (0.2679,-0.0187) -- (0.2679,0.0000) -- (0.2679,0.0187) -- (0.2676,0.0374) -- (0.2672,0.0562) -- (0.2665,0.0749) -- (0.2658,0.0937) -- (0.2648,0.1126) -- (0.2636,0.1315) -- (0.2623,0.1504) -- (0.2608,0.1695) -- (0.2590,0.1886) -- (0.2571,0.2078) -- (0.2550,0.2271) -- (0.2527,0.2465) -- (0.2502,0.2660) -- (0.2474,0.2857) -- (0.2445,0.3055) -- (0.2413,0.3255) -- (0.2379,0.3456) -- (0.2342,0.3659) -- (0.2303,0.3864) -- (0.2261,0.4071) -- (0.2216,0.4280) -- (0.2169,0.4491) -- (0.2118,0.4705) -- (0.2065,0.4921) -- (0.2008,0.5140) -- (0.1948,0.5361) -- (0.1884,0.5585) -- (0.1816,0.5813) -- (0.1745,0.6043) -- (0.1669,0.6277) -- (0.1589,0.6515) -- (0.1504,0.6756) -- (0.1414,0.7001) -- (0.1319,0.7250) -- (0.1219,0.7503) -- (0.1113,0.7760) -- (0.1000,0.8022) -- (0.0881,0.8289) -- (0.0755,0.8561) -- (0.0621,0.8838) -- (0.0479,0.9120) -- (0.0329,0.9407) -- (0.0169,0.9701) -- (0.0000,1.0000);
\draw[wulffA] (0.0000,-1.0000) -- (0.0116,-0.9677) -- (0.0224,-0.9362) -- (0.0326,-0.9056) -- (0.0421,-0.8757) -- (0.0511,-0.8467) -- (0.0595,-0.8183) -- (0.0674,-0.7906) -- (0.0749,-0.7635) -- (0.0819,-0.7370) -- (0.0885,-0.7111) -- (0.0948,-0.6858) -- (0.1006,-0.6609) -- (0.1062,-0.6366) -- (0.1114,-0.6127) -- (0.1163,-0.5892) -- (0.1210,-0.5661) -- (0.1254,-0.5435) -- (0.1295,-0.5212) -- (0.1334,-0.4992) -- (0.1371,-0.4776) -- (0.1405,-0.4563) -- (0.1438,-0.4352) -- (0.1468,-0.4145) -- (0.1497,-0.3940) -- (0.1523,-0.3737) -- (0.1549,-0.3537) -- (0.1572,-0.3339) -- (0.1594,-0.3143) -- (0.1614,-0.2949) -- (0.1633,-0.2757) -- (0.1650,-0.2566) -- (0.1667,-0.2377) -- (0.1681,-0.2189) -- (0.1695,-0.2002) -- (0.1707,-0.1816) -- (0.1718,-0.1632) -- (0.1727,-0.1448) -- (0.1736,-0.1265) -- (0.1743,-0.1083) -- (0.1749,-0.0902) -- (0.1754,-0.0721) -- (0.1758,-0.0540) -- (0.1761,-0.0360) -- (0.1763,-0.0180) -- (0.1763,0.0000) -- (0.1763,0.0180) -- (0.1761,0.0360) -- (0.1758,0.0540) -- (0.1754,0.0721) -- (0.1749,0.0902) -- (0.1743,0.1083) -- (0.1736,0.1265) -- (0.1727,0.1448) -- (0.1718,0.1632) -- (0.1707,0.1816) -- (0.1695,0.2002) -- (0.1681,0.2189) -- (0.1667,0.2377) -- (0.1650,0.2566) -- (0.1633,0.2757) -- (0.1614,0.2949) -- (0.1594,0.3143) -- (0.1572,0.3339) -- (0.1549,0.3537) -- (0.1523,0.3737) -- (0.1497,0.3940) -- (0.1468,0.4145) -- (0.1438,0.4352) -- (0.1405,0.4563) -- (0.1371,0.4776) -- (0.1334,0.4992) -- (0.1295,0.5212) -- (0.1254,0.5435) -- (0.1210,0.5661) -- (0.1163,0.5892) -- (0.1114,0.6127) -- (0.1062,0.6366) -- (0.1006,0.6609) -- (0.0948,0.6858) -- (0.0885,0.7111) -- (0.0819,0.7370) -- (0.0749,0.7635) -- (0.0674,0.7906) -- (0.0595,0.8183) -- (0.0511,0.8467) -- (0.0421,0.8757) -- (0.0326,0.9056) -- (0.0224,0.9362) -- (0.0116,0.9677) -- (0.0000,1.0000);
\draw[wulffA] (0.0000,-1.0000) -- (0.0059,-0.9662) -- (0.0113,-0.9334) -- (0.0165,-0.9017) -- (0.0213,-0.8709) -- (0.0258,-0.8410) -- (0.0300,-0.8119) -- (0.0339,-0.7836) -- (0.0376,-0.7560) -- (0.0411,-0.7292) -- (0.0444,-0.7029) -- (0.0475,-0.6773) -- (0.0504,-0.6523) -- (0.0532,-0.6278) -- (0.0557,-0.6038) -- (0.0582,-0.5803) -- (0.0605,-0.5572) -- (0.0626,-0.5346) -- (0.0646,-0.5124) -- (0.0666,-0.4906) -- (0.0684,-0.4691) -- (0.0700,-0.4480) -- (0.0716,-0.4271) -- (0.0731,-0.4066) -- (0.0745,-0.3864) -- (0.0758,-0.3664) -- (0.0770,-0.3466) -- (0.0782,-0.3271) -- (0.0793,-0.3079) -- (0.0802,-0.2888) -- (0.0812,-0.2699) -- (0.0820,-0.2511) -- (0.0828,-0.2325) -- (0.0835,-0.2141) -- (0.0842,-0.1958) -- (0.0847,-0.1776) -- (0.0853,-0.1596) -- (0.0857,-0.1416) -- (0.0862,-0.1237) -- (0.0865,-0.1059) -- (0.0868,-0.0882) -- (0.0871,-0.0705) -- (0.0872,-0.0528) -- (0.0874,-0.0352) -- (0.0875,-0.0176) -- (0.0875,0.0000) -- (0.0875,0.0176) -- (0.0874,0.0352) -- (0.0872,0.0528) -- (0.0871,0.0705) -- (0.0868,0.0882) -- (0.0865,0.1059) -- (0.0862,0.1237) -- (0.0857,0.1416) -- (0.0853,0.1596) -- (0.0847,0.1776) -- (0.0842,0.1958) -- (0.0835,0.2141) -- (0.0828,0.2325) -- (0.0820,0.2511) -- (0.0812,0.2699) -- (0.0802,0.2888) -- (0.0793,0.3079) -- (0.0782,0.3271) -- (0.0770,0.3466) -- (0.0758,0.3664) -- (0.0745,0.3864) -- (0.0731,0.4066) -- (0.0716,0.4271) -- (0.0700,0.4480) -- (0.0684,0.4691) -- (0.0666,0.4906) -- (0.0646,0.5124) -- (0.0626,0.5346) -- (0.0605,0.5572) -- (0.0582,0.5803) -- (0.0557,0.6038) -- (0.0532,0.6278) -- (0.0504,0.6523) -- (0.0475,0.6773) -- (0.0444,0.7029) -- (0.0411,0.7292) -- (0.0376,0.7560) -- (0.0339,0.7836) -- (0.0300,0.8119) -- (0.0258,0.8410) -- (0.0213,0.8709) -- (0.0165,0.9017) -- (0.0113,0.9334) -- (0.0059,0.9662) -- (0.0000,1.0000);
\draw[wulffB] (0.0000,-1.0000) -- (0.0000,-0.9657) -- (0.0000,-0.9325) -- (0.0000,-0.9004) -- (0.0000,-0.8693) -- (0.0000,-0.8391) -- (0.0000,-0.8098) -- (0.0000,-0.7813) -- (0.0000,-0.7536) -- (0.0000,-0.7265) -- (0.0000,-0.7002) -- (0.0000,-0.6745) -- (0.0000,-0.6494) -- (0.0000,-0.6249) -- (0.0000,-0.6009) -- (0.0000,-0.5774) -- (0.0000,-0.5543) -- (0.0000,-0.5317) -- (0.0000,-0.5095) -- (0.0000,-0.4877) -- (0.0000,-0.4663) -- (0.0000,-0.4452) -- (0.0000,-0.4245) -- (0.0000,-0.4040) -- (0.0000,-0.3839) -- (0.0000,-0.3640) -- (0.0000,-0.3443) -- (0.0000,-0.3249) -- (0.0000,-0.3057) -- (0.0000,-0.2867) -- (0.0000,-0.2679) -- (0.0000,-0.2493) -- (0.0000,-0.2309) -- (0.0000,-0.2126) -- (0.0000,-0.1944) -- (0.0000,-0.1763) -- (0.0000,-0.1584) -- (0.0000,-0.1405) -- (0.0000,-0.1228) -- (0.0000,-0.1051) -- (0.0000,-0.0875) -- (0.0000,-0.0699) -- (0.0000,-0.0524) -- (0.0000,-0.0349) -- (0.0000,-0.0175) -- (0.0000,0.0000) -- (0.0000,0.0175) -- (0.0000,0.0349) -- (0.0000,0.0524) -- (0.0000,0.0699) -- (0.0000,0.0875) -- (0.0000,0.1051) -- (0.0000,0.1228) -- (0.0000,0.1405) -- (0.0000,0.1584) -- (0.0000,0.1763) -- (0.0000,0.1944) -- (0.0000,0.2126) -- (0.0000,0.2309) -- (0.0000,0.2493) -- (0.0000,0.2679) -- (0.0000,0.2867) -- (0.0000,0.3057) -- (0.0000,0.3249) -- (0.0000,0.3443) -- (0.0000,0.3640) -- (0.0000,0.3839) -- (0.0000,0.4040) -- (0.0000,0.4245) -- (0.0000,0.4452) -- (0.0000,0.4663) -- (0.0000,0.4877) -- (0.0000,0.5095) -- (0.0000,0.5317) -- (0.0000,0.5543) -- (0.0000,0.5774) -- (0.0000,0.6009) -- (0.0000,0.6249) -- (0.0000,0.6494) -- (0.0000,0.6745) -- (0.0000,0.7002) -- (0.0000,0.7265) -- (0.0000,0.7536) -- (0.0000,0.7813) -- (0.0000,0.8098) -- (0.0000,0.8391) -- (0.0000,0.8693) -- (0.0000,0.9004) -- (0.0000,0.9325) -- (0.0000,0.9657) -- (0.0000,1.0000);
\draw[wulffA] (-0.0000,-1.0000) -- (-0.0059,-0.9662) -- (-0.0113,-0.9334) -- (-0.0165,-0.9017) -- (-0.0213,-0.8709) -- (-0.0258,-0.8410) -- (-0.0300,-0.8119) -- (-0.0339,-0.7836) -- (-0.0376,-0.7560) -- (-0.0411,-0.7292) -- (-0.0444,-0.7029) -- (-0.0475,-0.6773) -- (-0.0504,-0.6523) -- (-0.0532,-0.6278) -- (-0.0557,-0.6038) -- (-0.0582,-0.5803) -- (-0.0605,-0.5572) -- (-0.0626,-0.5346) -- (-0.0646,-0.5124) -- (-0.0666,-0.4906) -- (-0.0684,-0.4691) -- (-0.0700,-0.4480) -- (-0.0716,-0.4271) -- (-0.0731,-0.4066) -- (-0.0745,-0.3864) -- (-0.0758,-0.3664) -- (-0.0770,-0.3466) -- (-0.0782,-0.3271) -- (-0.0793,-0.3079) -- (-0.0802,-0.2888) -- (-0.0812,-0.2699) -- (-0.0820,-0.2511) -- (-0.0828,-0.2325) -- (-0.0835,-0.2141) -- (-0.0842,-0.1958) -- (-0.0847,-0.1776) -- (-0.0853,-0.1596) -- (-0.0857,-0.1416) -- (-0.0862,-0.1237) -- (-0.0865,-0.1059) -- (-0.0868,-0.0882) -- (-0.0871,-0.0705) -- (-0.0872,-0.0528) -- (-0.0874,-0.0352) -- (-0.0875,-0.0176) -- (-0.0875,0.0000) -- (-0.0875,0.0176) -- (-0.0874,0.0352) -- (-0.0872,0.0528) -- (-0.0871,0.0705) -- (-0.0868,0.0882) -- (-0.0865,0.1059) -- (-0.0862,0.1237) -- (-0.0857,0.1416) -- (-0.0853,0.1596) -- (-0.0847,0.1776) -- (-0.0842,0.1958) -- (-0.0835,0.2141) -- (-0.0828,0.2325) -- (-0.0820,0.2511) -- (-0.0812,0.2699) -- (-0.0802,0.2888) -- (-0.0793,0.3079) -- (-0.0782,0.3271) -- (-0.0770,0.3466) -- (-0.0758,0.3664) -- (-0.0745,0.3864) -- (-0.0731,0.4066) -- (-0.0716,0.4271) -- (-0.0700,0.4480) -- (-0.0684,0.4691) -- (-0.0666,0.4906) -- (-0.0646,0.5124) -- (-0.0626,0.5346) -- (-0.0605,0.5572) -- (-0.0582,0.5803) -- (-0.0557,0.6038) -- (-0.0532,0.6278) -- (-0.0504,0.6523) -- (-0.0475,0.6773) -- (-0.0444,0.7029) -- (-0.0411,0.7292) -- (-0.0376,0.7560) -- (-0.0339,0.7836) -- (-0.0300,0.8119) -- (-0.0258,0.8410) -- (-0.0213,0.8709) -- (-0.0165,0.9017) -- (-0.0113,0.9334) -- (-0.0059,0.9662) -- (-0.0000,1.0000);
\draw[wulffA] (-0.0000,-1.0000) -- (-0.0116,-0.9677) -- (-0.0224,-0.9362) -- (-0.0326,-0.9056) -- (-0.0421,-0.8757) -- (-0.0511,-0.8467) -- (-0.0595,-0.8183) -- (-0.0674,-0.7906) -- (-0.0749,-0.7635) -- (-0.0819,-0.7370) -- (-0.0885,-0.7111) -- (-0.0948,-0.6858) -- (-0.1006,-0.6609) -- (-0.1062,-0.6366) -- (-0.1114,-0.6127) -- (-0.1163,-0.5892) -- (-0.1210,-0.5661) -- (-0.1254,-0.5435) -- (-0.1295,-0.5212) -- (-0.1334,-0.4992) -- (-0.1371,-0.4776) -- (-0.1405,-0.4563) -- (-0.1438,-0.4352) -- (-0.1468,-0.4145) -- (-0.1497,-0.3940) -- (-0.1523,-0.3737) -- (-0.1549,-0.3537) -- (-0.1572,-0.3339) -- (-0.1594,-0.3143) -- (-0.1614,-0.2949) -- (-0.1633,-0.2757) -- (-0.1650,-0.2566) -- (-0.1667,-0.2377) -- (-0.1681,-0.2189) -- (-0.1695,-0.2002) -- (-0.1707,-0.1816) -- (-0.1718,-0.1632) -- (-0.1727,-0.1448) -- (-0.1736,-0.1265) -- (-0.1743,-0.1083) -- (-0.1749,-0.0902) -- (-0.1754,-0.0721) -- (-0.1758,-0.0540) -- (-0.1761,-0.0360) -- (-0.1763,-0.0180) -- (-0.1763,0.0000) -- (-0.1763,0.0180) -- (-0.1761,0.0360) -- (-0.1758,0.0540) -- (-0.1754,0.0721) -- (-0.1749,0.0902) -- (-0.1743,0.1083) -- (-0.1736,0.1265) -- (-0.1727,0.1448) -- (-0.1718,0.1632) -- (-0.1707,0.1816) -- (-0.1695,0.2002) -- (-0.1681,0.2189) -- (-0.1667,0.2377) -- (-0.1650,0.2566) -- (-0.1633,0.2757) -- (-0.1614,0.2949) -- (-0.1594,0.3143) -- (-0.1572,0.3339) -- (-0.1549,0.3537) -- (-0.1523,0.3737) -- (-0.1497,0.3940) -- (-0.1468,0.4145) -- (-0.1438,0.4352) -- (-0.1405,0.4563) -- (-0.1371,0.4776) -- (-0.1334,0.4992) -- (-0.1295,0.5212) -- (-0.1254,0.5435) -- (-0.1210,0.5661) -- (-0.1163,0.5892) -- (-0.1114,0.6127) -- (-0.1062,0.6366) -- (-0.1006,0.6609) -- (-0.0948,0.6858) -- (-0.0885,0.7111) -- (-0.0819,0.7370) -- (-0.0749,0.7635) -- (-0.0674,0.7906) -- (-0.0595,0.8183) -- (-0.0511,0.8467) -- (-0.0421,0.8757) -- (-0.0326,0.9056) -- (-0.0224,0.9362) -- (-0.0116,0.9677) -- (-0.0000,1.0000);
\draw[wulffB] (-0.0000,-1.0000) -- (-0.0169,-0.9701) -- (-0.0329,-0.9407) -- (-0.0479,-0.9120) -- (-0.0621,-0.8838) -- (-0.0755,-0.8561) -- (-0.0881,-0.8289) -- (-0.1000,-0.8022) -- (-0.1113,-0.7760) -- (-0.1219,-0.7503) -- (-0.1319,-0.7250) -- (-0.1414,-0.7001) -- (-0.1504,-0.6756) -- (-0.1589,-0.6515) -- (-0.1669,-0.6277) -- (-0.1745,-0.6043) -- (-0.1816,-0.5813) -- (-0.1884,-0.5585) -- (-0.1948,-0.5361) -- (-0.2008,-0.5140) -- (-0.2065,-0.4921) -- (-0.2118,-0.4705) -- (-0.2169,-0.4491) -- (-0.2216,-0.4280) -- (-0.2261,-0.4071) -- (-0.2303,-0.3864) -- (-0.2342,-0.3659) -- (-0.2379,-0.3456) -- (-0.2413,-0.3255) -- (-0.2445,-0.3055) -- (-0.2474,-0.2857) -- (-0.2502,-0.2660) -- (-0.2527,-0.2465) -- (-0.2550,-0.2271) -- (-0.2571,-0.2078) -- (-0.2590,-0.1886) -- (-0.2608,-0.1695) -- (-0.2623,-0.1504) -- (-0.2636,-0.1315) -- (-0.2648,-0.1126) -- (-0.2658,-0.0937) -- (-0.2665,-0.0749) -- (-0.2672,-0.0562) -- (-0.2676,-0.0374) -- (-0.2679,-0.0187) -- (-0.2679,0.0000) -- (-0.2679,0.0187) -- (-0.2676,0.0374) -- (-0.2672,0.0562) -- (-0.2665,0.0749) -- (-0.2658,0.0937) -- (-0.2648,0.1126) -- (-0.2636,0.1315) -- (-0.2623,0.1504) -- (-0.2608,0.1695) -- (-0.2590,0.1886) -- (-0.2571,0.2078) -- (-0.2550,0.2271) -- (-0.2527,0.2465) -- (-0.2502,0.2660) -- (-0.2474,0.2857) -- (-0.2445,0.3055) -- (-0.2413,0.3255) -- (-0.2379,0.3456) -- (-0.2342,0.3659) -- (-0.2303,0.3864) -- (-0.2261,0.4071) -- (-0.2216,0.4280) -- (-0.2169,0.4491) -- (-0.2118,0.4705) -- (-0.2065,0.4921) -- (-0.2008,0.5140) -- (-0.1948,0.5361) -- (-0.1884,0.5585) -- (-0.1816,0.5813) -- (-0.1745,0.6043) -- (-0.1669,0.6277) -- (-0.1589,0.6515) -- (-0.1504,0.6756) -- (-0.1414,0.7001) -- (-0.1319,0.7250) -- (-0.1219,0.7503) -- (-0.1113,0.7760) -- (-0.1000,0.8022) -- (-0.0881,0.8289) -- (-0.0755,0.8561) -- (-0.0621,0.8838) -- (-0.0479,0.9120) -- (-0.0329,0.9407) -- (-0.0169,0.9701) -- (-0.0000,1.0000);
\draw[wulffA] (-0.0000,-1.0000) -- (-0.0218,-0.9734) -- (-0.0426,-0.9470) -- (-0.0622,-0.9208) -- (-0.0808,-0.8949) -- (-0.0985,-0.8692) -- (-0.1153,-0.8438) -- (-0.1312,-0.8186) -- (-0.1463,-0.7937) -- (-0.1606,-0.7690) -- (-0.1742,-0.7446) -- (-0.1871,-0.7204) -- (-0.1993,-0.6965) -- (-0.2109,-0.6728) -- (-0.2219,-0.6494) -- (-0.2324,-0.6262) -- (-0.2423,-0.6032) -- (-0.2516,-0.5804) -- (-0.2605,-0.5578) -- (-0.2689,-0.5355) -- (-0.2769,-0.5133) -- (-0.2844,-0.4913) -- (-0.2914,-0.4695) -- (-0.2981,-0.4479) -- (-0.3044,-0.4264) -- (-0.3103,-0.4051) -- (-0.3159,-0.3839) -- (-0.3211,-0.3629) -- (-0.3259,-0.3420) -- (-0.3304,-0.3212) -- (-0.3347,-0.3006) -- (-0.3386,-0.2801) -- (-0.3422,-0.2596) -- (-0.3455,-0.2393) -- (-0.3485,-0.2190) -- (-0.3512,-0.1989) -- (-0.3537,-0.1788) -- (-0.3559,-0.1587) -- (-0.3578,-0.1388) -- (-0.3594,-0.1189) -- (-0.3608,-0.0990) -- (-0.3620,-0.0791) -- (-0.3628,-0.0593) -- (-0.3635,-0.0395) -- (-0.3638,-0.0198) -- (-0.3640,0.0000) -- (-0.3638,0.0198) -- (-0.3635,0.0395) -- (-0.3628,0.0593) -- (-0.3620,0.0791) -- (-0.3608,0.0990) -- (-0.3594,0.1189) -- (-0.3578,0.1388) -- (-0.3559,0.1587) -- (-0.3537,0.1788) -- (-0.3512,0.1989) -- (-0.3485,0.2190) -- (-0.3455,0.2393) -- (-0.3422,0.2596) -- (-0.3386,0.2801) -- (-0.3347,0.3006) -- (-0.3304,0.3212) -- (-0.3259,0.3420) -- (-0.3211,0.3629) -- (-0.3159,0.3839) -- (-0.3103,0.4051) -- (-0.3044,0.4264) -- (-0.2981,0.4479) -- (-0.2914,0.4695) -- (-0.2844,0.4913) -- (-0.2769,0.5133) -- (-0.2689,0.5355) -- (-0.2605,0.5578) -- (-0.2516,0.5804) -- (-0.2423,0.6032) -- (-0.2324,0.6262) -- (-0.2219,0.6494) -- (-0.2109,0.6728) -- (-0.1993,0.6965) -- (-0.1871,0.7204) -- (-0.1742,0.7446) -- (-0.1606,0.7690) -- (-0.1463,0.7937) -- (-0.1312,0.8186) -- (-0.1153,0.8438) -- (-0.0985,0.8692) -- (-0.0808,0.8949) -- (-0.0622,0.9208) -- (-0.0426,0.9470) -- (-0.0218,0.9734) -- (-0.0000,1.0000);
\draw[wulffA] (-0.0000,-1.0000) -- (-0.0261,-0.9775) -- (-0.0511,-0.9548) -- (-0.0750,-0.9319) -- (-0.0979,-0.9090) -- (-0.1197,-0.8859) -- (-0.1405,-0.8628) -- (-0.1604,-0.8397) -- (-0.1794,-0.8166) -- (-0.1975,-0.7935) -- (-0.2148,-0.7703) -- (-0.2313,-0.7473) -- (-0.2470,-0.7242) -- (-0.2620,-0.7012) -- (-0.2763,-0.6783) -- (-0.2899,-0.6554) -- (-0.3028,-0.6326) -- (-0.3151,-0.6098) -- (-0.3268,-0.5872) -- (-0.3379,-0.5646) -- (-0.3484,-0.5421) -- (-0.3584,-0.5196) -- (-0.3679,-0.4973) -- (-0.3768,-0.4750) -- (-0.3853,-0.4528) -- (-0.3932,-0.4307) -- (-0.4007,-0.4087) -- (-0.4077,-0.3867) -- (-0.4143,-0.3648) -- (-0.4204,-0.3430) -- (-0.4262,-0.3212) -- (-0.4315,-0.2995) -- (-0.4364,-0.2778) -- (-0.4409,-0.2563) -- (-0.4450,-0.2347) -- (-0.4488,-0.2132) -- (-0.4521,-0.1918) -- (-0.4551,-0.1704) -- (-0.4578,-0.1490) -- (-0.4601,-0.1277) -- (-0.4620,-0.1063) -- (-0.4635,-0.0850) -- (-0.4647,-0.0638) -- (-0.4656,-0.0425) -- (-0.4661,-0.0212) -- (-0.4663,0.0000) -- (-0.4661,0.0212) -- (-0.4656,0.0425) -- (-0.4647,0.0638) -- (-0.4635,0.0850) -- (-0.4620,0.1063) -- (-0.4601,0.1277) -- (-0.4578,0.1490) -- (-0.4551,0.1704) -- (-0.4521,0.1918) -- (-0.4488,0.2132) -- (-0.4450,0.2347) -- (-0.4409,0.2563) -- (-0.4364,0.2778) -- (-0.4315,0.2995) -- (-0.4262,0.3212) -- (-0.4204,0.3430) -- (-0.4143,0.3648) -- (-0.4077,0.3867) -- (-0.4007,0.4087) -- (-0.3932,0.4307) -- (-0.3853,0.4528) -- (-0.3768,0.4750) -- (-0.3679,0.4973) -- (-0.3584,0.5196) -- (-0.3484,0.5421) -- (-0.3379,0.5646) -- (-0.3268,0.5872) -- (-0.3151,0.6098) -- (-0.3028,0.6326) -- (-0.2899,0.6554) -- (-0.2763,0.6783) -- (-0.2620,0.7012) -- (-0.2470,0.7242) -- (-0.2313,0.7473) -- (-0.2148,0.7703) -- (-0.1975,0.7935) -- (-0.1794,0.8166) -- (-0.1604,0.8397) -- (-0.1405,0.8628) -- (-0.1197,0.8859) -- (-0.0979,0.9090) -- (-0.0750,0.9319) -- (-0.0511,0.9548) -- (-0.0261,0.9775) -- (-0.0000,1.0000);
\draw[wulffB] (-0.0000,-1.0000) -- (-0.0297,-0.9823) -- (-0.0584,-0.9639) -- (-0.0860,-0.9451) -- (-0.1127,-0.9258) -- (-0.1384,-0.9061) -- (-0.1631,-0.8860) -- (-0.1869,-0.8656) -- (-0.2098,-0.8448) -- (-0.2318,-0.8238) -- (-0.2529,-0.8025) -- (-0.2732,-0.7809) -- (-0.2927,-0.7592) -- (-0.3114,-0.7372) -- (-0.3293,-0.7151) -- (-0.3464,-0.6928) -- (-0.3628,-0.6704) -- (-0.3785,-0.6479) -- (-0.3934,-0.6253) -- (-0.4077,-0.6025) -- (-0.4213,-0.5797) -- (-0.4342,-0.5568) -- (-0.4465,-0.5339) -- (-0.4582,-0.5109) -- (-0.4692,-0.4879) -- (-0.4797,-0.4648) -- (-0.4896,-0.4416) -- (-0.4988,-0.4185) -- (-0.5076,-0.3953) -- (-0.5157,-0.3721) -- (-0.5234,-0.3489) -- (-0.5305,-0.3257) -- (-0.5370,-0.3025) -- (-0.5431,-0.2792) -- (-0.5486,-0.2560) -- (-0.5537,-0.2327) -- (-0.5582,-0.2094) -- (-0.5622,-0.1862) -- (-0.5658,-0.1629) -- (-0.5689,-0.1396) -- (-0.5715,-0.1164) -- (-0.5736,-0.0931) -- (-0.5752,-0.0698) -- (-0.5764,-0.0465) -- (-0.5771,-0.0233) -- (-0.5774,0.0000) -- (-0.5771,0.0233) -- (-0.5764,0.0465) -- (-0.5752,0.0698) -- (-0.5736,0.0931) -- (-0.5715,0.1164) -- (-0.5689,0.1396) -- (-0.5658,0.1629) -- (-0.5622,0.1862) -- (-0.5582,0.2094) -- (-0.5537,0.2327) -- (-0.5486,0.2560) -- (-0.5431,0.2792) -- (-0.5370,0.3025) -- (-0.5305,0.3257) -- (-0.5234,0.3489) -- (-0.5157,0.3721) -- (-0.5076,0.3953) -- (-0.4988,0.4185) -- (-0.4896,0.4416) -- (-0.4797,0.4648) -- (-0.4692,0.4879) -- (-0.4582,0.5109) -- (-0.4465,0.5339) -- (-0.4342,0.5568) -- (-0.4213,0.5797) -- (-0.4077,0.6025) -- (-0.3934,0.6253) -- (-0.3785,0.6479) -- (-0.3628,0.6704) -- (-0.3464,0.6928) -- (-0.3293,0.7151) -- (-0.3114,0.7372) -- (-0.2927,0.7592) -- (-0.2732,0.7809) -- (-0.2529,0.8025) -- (-0.2318,0.8238) -- (-0.2098,0.8448) -- (-0.1869,0.8656) -- (-0.1631,0.8860) -- (-0.1384,0.9061) -- (-0.1127,0.9258) -- (-0.0860,0.9451) -- (-0.0584,0.9639) -- (-0.0297,0.9823) -- (-0.0000,1.0000);
\draw[wulffA] (-0.0000,-1.0000) -- (-0.0324,-0.9876) -- (-0.0640,-0.9743) -- (-0.0948,-0.9602) -- (-0.1248,-0.9453) -- (-0.1540,-0.9296) -- (-0.1824,-0.9132) -- (-0.2100,-0.8961) -- (-0.2367,-0.8784) -- (-0.2626,-0.8601) -- (-0.2877,-0.8413) -- (-0.3120,-0.8219) -- (-0.3355,-0.8020) -- (-0.3582,-0.7816) -- (-0.3801,-0.7608) -- (-0.4012,-0.7396) -- (-0.4216,-0.7179) -- (-0.4411,-0.6959) -- (-0.4599,-0.6736) -- (-0.4779,-0.6509) -- (-0.4952,-0.6280) -- (-0.5117,-0.6047) -- (-0.5274,-0.5812) -- (-0.5425,-0.5575) -- (-0.5568,-0.5335) -- (-0.5704,-0.5093) -- (-0.5833,-0.4850) -- (-0.5955,-0.4604) -- (-0.6069,-0.4357) -- (-0.6177,-0.4108) -- (-0.6278,-0.3857) -- (-0.6373,-0.3606) -- (-0.6460,-0.3353) -- (-0.6541,-0.3099) -- (-0.6615,-0.2844) -- (-0.6683,-0.2588) -- (-0.6743,-0.2332) -- (-0.6798,-0.2074) -- (-0.6846,-0.1816) -- (-0.6887,-0.1558) -- (-0.6923,-0.1299) -- (-0.6951,-0.1040) -- (-0.6973,-0.0780) -- (-0.6989,-0.0520) -- (-0.6999,-0.0260) -- (-0.7002,0.0000) -- (-0.6999,0.0260) -- (-0.6989,0.0520) -- (-0.6973,0.0780) -- (-0.6951,0.1040) -- (-0.6923,0.1299) -- (-0.6887,0.1558) -- (-0.6846,0.1816) -- (-0.6798,0.2074) -- (-0.6743,0.2332) -- (-0.6683,0.2588) -- (-0.6615,0.2844) -- (-0.6541,0.3099) -- (-0.6460,0.3353) -- (-0.6373,0.3606) -- (-0.6278,0.3857) -- (-0.6177,0.4108) -- (-0.6069,0.4357) -- (-0.5955,0.4604) -- (-0.5833,0.4850) -- (-0.5704,0.5093) -- (-0.5568,0.5335) -- (-0.5425,0.5575) -- (-0.5274,0.5812) -- (-0.5117,0.6047) -- (-0.4952,0.6280) -- (-0.4779,0.6509) -- (-0.4599,0.6736) -- (-0.4411,0.6959) -- (-0.4216,0.7179) -- (-0.4012,0.7396) -- (-0.3801,0.7608) -- (-0.3582,0.7816) -- (-0.3355,0.8020) -- (-0.3120,0.8219) -- (-0.2877,0.8413) -- (-0.2626,0.8601) -- (-0.2367,0.8784) -- (-0.2100,0.8961) -- (-0.1824,0.9132) -- (-0.1540,0.9296) -- (-0.1248,0.9453) -- (-0.0948,0.9602) -- (-0.0640,0.9743) -- (-0.0324,0.9876) -- (-0.0000,1.0000);
\draw[wulffA] (-0.0000,-1.0000) -- (-0.0342,-0.9934) -- (-0.0679,-0.9856) -- (-0.1011,-0.9768) -- (-0.1338,-0.9669) -- (-0.1660,-0.9560) -- (-0.1976,-0.9441) -- (-0.2286,-0.9312) -- (-0.2591,-0.9174) -- (-0.2888,-0.9026) -- (-0.3179,-0.8870) -- (-0.3464,-0.8706) -- (-0.3741,-0.8533) -- (-0.4012,-0.8352) -- (-0.4275,-0.8164) -- (-0.4531,-0.7968) -- (-0.4779,-0.7766) -- (-0.5020,-0.7557) -- (-0.5252,-0.7341) -- (-0.5477,-0.7119) -- (-0.5695,-0.6891) -- (-0.5904,-0.6658) -- (-0.6105,-0.6419) -- (-0.6297,-0.6175) -- (-0.6482,-0.5927) -- (-0.6658,-0.5673) -- (-0.6826,-0.5416) -- (-0.6986,-0.5154) -- (-0.7137,-0.4888) -- (-0.7280,-0.4619) -- (-0.7414,-0.4346) -- (-0.7539,-0.4071) -- (-0.7656,-0.3792) -- (-0.7765,-0.3510) -- (-0.7865,-0.3227) -- (-0.7956,-0.2940) -- (-0.8039,-0.2652) -- (-0.8112,-0.2362) -- (-0.8178,-0.2070) -- (-0.8234,-0.1777) -- (-0.8282,-0.1483) -- (-0.8321,-0.1188) -- (-0.8352,-0.0891) -- (-0.8374,-0.0595) -- (-0.8387,-0.0297) -- (-0.8391,0.0000) -- (-0.8387,0.0297) -- (-0.8374,0.0595) -- (-0.8352,0.0891) -- (-0.8321,0.1188) -- (-0.8282,0.1483) -- (-0.8234,0.1777) -- (-0.8178,0.2070) -- (-0.8112,0.2362) -- (-0.8039,0.2652) -- (-0.7956,0.2940) -- (-0.7865,0.3227) -- (-0.7765,0.3510) -- (-0.7656,0.3792) -- (-0.7539,0.4071) -- (-0.7414,0.4346) -- (-0.7280,0.4619) -- (-0.7137,0.4888) -- (-0.6986,0.5154) -- (-0.6826,0.5416) -- (-0.6658,0.5673) -- (-0.6482,0.5927) -- (-0.6297,0.6175) -- (-0.6105,0.6419) -- (-0.5904,0.6658) -- (-0.5695,0.6891) -- (-0.5477,0.7119) -- (-0.5252,0.7341) -- (-0.5020,0.7557) -- (-0.4779,0.7766) -- (-0.4531,0.7968) -- (-0.4275,0.8164) -- (-0.4012,0.8352) -- (-0.3741,0.8533) -- (-0.3464,0.8706) -- (-0.3179,0.8870) -- (-0.2888,0.9026) -- (-0.2591,0.9174) -- (-0.2286,0.9312) -- (-0.1976,0.9441) -- (-0.1660,0.9560) -- (-0.1338,0.9669) -- (-0.1011,0.9768) -- (-0.0679,0.9856) -- (-0.0342,0.9934) -- (-0.0000,1.0000);
\draw[wulffA] (0.1736,0.9848) -- (0.1725,0.9789) -- (0.1712,0.9730) -- (0.1696,0.9673) -- (0.1679,0.9616) -- (0.1660,0.9560) -- (0.1639,0.9505) -- (0.1617,0.9451) -- (0.1593,0.9398) -- (0.1567,0.9347) -- (0.1540,0.9296) -- (0.1512,0.9247) -- (0.1482,0.9198) -- (0.1450,0.9151) -- (0.1418,0.9106) -- (0.1384,0.9061) -- (0.1349,0.9018) -- (0.1312,0.8976) -- (0.1275,0.8936) -- (0.1236,0.8897) -- (0.1197,0.8859) -- (0.1156,0.8823) -- (0.1115,0.8788) -- (0.1072,0.8755) -- (0.1029,0.8722) -- (0.0985,0.8692) -- (0.0940,0.8663) -- (0.0895,0.8635) -- (0.0849,0.8609) -- (0.0802,0.8584) -- (0.0755,0.8561) -- (0.0707,0.8539) -- (0.0658,0.8519) -- (0.0610,0.8500) -- (0.0560,0.8482) -- (0.0511,0.8467) -- (0.0461,0.8452) -- (0.0410,0.8439) -- (0.0360,0.8428) -- (0.0309,0.8418) -- (0.0258,0.8410) -- (0.0206,0.8403) -- (0.0155,0.8398) -- (0.0103,0.8394) -- (0.0052,0.8392) -- (0.0000,0.8391) -- (-0.0052,0.8392) -- (-0.0103,0.8394) -- (-0.0155,0.8398) -- (-0.0206,0.8403) -- (-0.0258,0.8410) -- (-0.0309,0.8418) -- (-0.0360,0.8428) -- (-0.0410,0.8439) -- (-0.0461,0.8452) -- (-0.0511,0.8467) -- (-0.0560,0.8482) -- (-0.0610,0.8500) -- (-0.0658,0.8519) -- (-0.0707,0.8539) -- (-0.0755,0.8561) -- (-0.0802,0.8584) -- (-0.0849,0.8609) -- (-0.0895,0.8635) -- (-0.0940,0.8663) -- (-0.0985,0.8692) -- (-0.1029,0.8722) -- (-0.1072,0.8755) -- (-0.1115,0.8788) -- (-0.1156,0.8823) -- (-0.1197,0.8859) -- (-0.1236,0.8897) -- (-0.1275,0.8936) -- (-0.1312,0.8976) -- (-0.1349,0.9018) -- (-0.1384,0.9061) -- (-0.1418,0.9106) -- (-0.1450,0.9151) -- (-0.1482,0.9198) -- (-0.1512,0.9247) -- (-0.1540,0.9296) -- (-0.1567,0.9347) -- (-0.1593,0.9398) -- (-0.1617,0.9451) -- (-0.1639,0.9505) -- (-0.1660,0.9560) -- (-0.1679,0.9616) -- (-0.1696,0.9673) -- (-0.1712,0.9730) -- (-0.1725,0.9789) -- (-0.1736,0.9848);
\draw[wulffA] (0.1736,-0.9848) -- (0.1725,-0.9789) -- (0.1712,-0.9730) -- (0.1696,-0.9673) -- (0.1679,-0.9616) -- (0.1660,-0.9560) -- (0.1639,-0.9505) -- (0.1617,-0.9451) -- (0.1593,-0.9398) -- (0.1567,-0.9347) -- (0.1540,-0.9296) -- (0.1512,-0.9247) -- (0.1482,-0.9198) -- (0.1450,-0.9151) -- (0.1418,-0.9106) -- (0.1384,-0.9061) -- (0.1349,-0.9018) -- (0.1312,-0.8976) -- (0.1275,-0.8936) -- (0.1236,-0.8897) -- (0.1197,-0.8859) -- (0.1156,-0.8823) -- (0.1115,-0.8788) -- (0.1072,-0.8755) -- (0.1029,-0.8722) -- (0.0985,-0.8692) -- (0.0940,-0.8663) -- (0.0895,-0.8635) -- (0.0849,-0.8609) -- (0.0802,-0.8584) -- (0.0755,-0.8561) -- (0.0707,-0.8539) -- (0.0658,-0.8519) -- (0.0610,-0.8500) -- (0.0560,-0.8482) -- (0.0511,-0.8467) -- (0.0461,-0.8452) -- (0.0410,-0.8439) -- (0.0360,-0.8428) -- (0.0309,-0.8418) -- (0.0258,-0.8410) -- (0.0206,-0.8403) -- (0.0155,-0.8398) -- (0.0103,-0.8394) -- (0.0052,-0.8392) -- (0.0000,-0.8391) -- (-0.0052,-0.8392) -- (-0.0103,-0.8394) -- (-0.0155,-0.8398) -- (-0.0206,-0.8403) -- (-0.0258,-0.8410) -- (-0.0309,-0.8418) -- (-0.0360,-0.8428) -- (-0.0410,-0.8439) -- (-0.0461,-0.8452) -- (-0.0511,-0.8467) -- (-0.0560,-0.8482) -- (-0.0610,-0.8500) -- (-0.0658,-0.8519) -- (-0.0707,-0.8539) -- (-0.0755,-0.8561) -- (-0.0802,-0.8584) -- (-0.0849,-0.8609) -- (-0.0895,-0.8635) -- (-0.0940,-0.8663) -- (-0.0985,-0.8692) -- (-0.1029,-0.8722) -- (-0.1072,-0.8755) -- (-0.1115,-0.8788) -- (-0.1156,-0.8823) -- (-0.1197,-0.8859) -- (-0.1236,-0.8897) -- (-0.1275,-0.8936) -- (-0.1312,-0.8976) -- (-0.1349,-0.9018) -- (-0.1384,-0.9061) -- (-0.1418,-0.9106) -- (-0.1450,-0.9151) -- (-0.1482,-0.9198) -- (-0.1512,-0.9247) -- (-0.1540,-0.9296) -- (-0.1567,-0.9347) -- (-0.1593,-0.9398) -- (-0.1617,-0.9451) -- (-0.1639,-0.9505) -- (-0.1660,-0.9560) -- (-0.1679,-0.9616) -- (-0.1696,-0.9673) -- (-0.1712,-0.9730) -- (-0.1725,-0.9789) -- (-0.1736,-0.9848);
\draw[wulffA] (0.3420,0.9397) -- (0.3378,0.9286) -- (0.3332,0.9178) -- (0.3284,0.9073) -- (0.3233,0.8970) -- (0.3179,0.8870) -- (0.3123,0.8773) -- (0.3065,0.8679) -- (0.3004,0.8587) -- (0.2942,0.8499) -- (0.2877,0.8413) -- (0.2811,0.8330) -- (0.2743,0.8249) -- (0.2673,0.8172) -- (0.2602,0.8097) -- (0.2529,0.8025) -- (0.2455,0.7955) -- (0.2380,0.7888) -- (0.2304,0.7824) -- (0.2226,0.7762) -- (0.2148,0.7703) -- (0.2068,0.7647) -- (0.1988,0.7593) -- (0.1907,0.7542) -- (0.1825,0.7493) -- (0.1742,0.7446) -- (0.1659,0.7402) -- (0.1575,0.7360) -- (0.1490,0.7321) -- (0.1405,0.7284) -- (0.1319,0.7250) -- (0.1233,0.7217) -- (0.1147,0.7187) -- (0.1060,0.7160) -- (0.0973,0.7134) -- (0.0885,0.7111) -- (0.0797,0.7091) -- (0.0709,0.7072) -- (0.0621,0.7055) -- (0.0533,0.7041) -- (0.0444,0.7029) -- (0.0356,0.7019) -- (0.0267,0.7012) -- (0.0178,0.7006) -- (0.0089,0.7003) -- (0.0000,0.7002) -- (-0.0089,0.7003) -- (-0.0178,0.7006) -- (-0.0267,0.7012) -- (-0.0356,0.7019) -- (-0.0444,0.7029) -- (-0.0533,0.7041) -- (-0.0621,0.7055) -- (-0.0709,0.7072) -- (-0.0797,0.7091) -- (-0.0885,0.7111) -- (-0.0973,0.7134) -- (-0.1060,0.7160) -- (-0.1147,0.7187) -- (-0.1233,0.7217) -- (-0.1319,0.7250) -- (-0.1405,0.7284) -- (-0.1490,0.7321) -- (-0.1575,0.7360) -- (-0.1659,0.7402) -- (-0.1742,0.7446) -- (-0.1825,0.7493) -- (-0.1907,0.7542) -- (-0.1988,0.7593) -- (-0.2068,0.7647) -- (-0.2148,0.7703) -- (-0.2226,0.7762) -- (-0.2304,0.7824) -- (-0.2380,0.7888) -- (-0.2455,0.7955) -- (-0.2529,0.8025) -- (-0.2602,0.8097) -- (-0.2673,0.8172) -- (-0.2743,0.8249) -- (-0.2811,0.8330) -- (-0.2877,0.8413) -- (-0.2942,0.8499) -- (-0.3004,0.8587) -- (-0.3065,0.8679) -- (-0.3123,0.8773) -- (-0.3179,0.8870) -- (-0.3233,0.8970) -- (-0.3284,0.9073) -- (-0.3332,0.9178) -- (-0.3378,0.9286) -- (-0.3420,0.9397);
\draw[wulffA] (0.3420,-0.9397) -- (0.3378,-0.9286) -- (0.3332,-0.9178) -- (0.3284,-0.9073) -- (0.3233,-0.8970) -- (0.3179,-0.8870) -- (0.3123,-0.8773) -- (0.3065,-0.8679) -- (0.3004,-0.8587) -- (0.2942,-0.8499) -- (0.2877,-0.8413) -- (0.2811,-0.8330) -- (0.2743,-0.8249) -- (0.2673,-0.8172) -- (0.2602,-0.8097) -- (0.2529,-0.8025) -- (0.2455,-0.7955) -- (0.2380,-0.7888) -- (0.2304,-0.7824) -- (0.2226,-0.7762) -- (0.2148,-0.7703) -- (0.2068,-0.7647) -- (0.1988,-0.7593) -- (0.1907,-0.7542) -- (0.1825,-0.7493) -- (0.1742,-0.7446) -- (0.1659,-0.7402) -- (0.1575,-0.7360) -- (0.1490,-0.7321) -- (0.1405,-0.7284) -- (0.1319,-0.7250) -- (0.1233,-0.7217) -- (0.1147,-0.7187) -- (0.1060,-0.7160) -- (0.0973,-0.7134) -- (0.0885,-0.7111) -- (0.0797,-0.7091) -- (0.0709,-0.7072) -- (0.0621,-0.7055) -- (0.0533,-0.7041) -- (0.0444,-0.7029) -- (0.0356,-0.7019) -- (0.0267,-0.7012) -- (0.0178,-0.7006) -- (0.0089,-0.7003) -- (0.0000,-0.7002) -- (-0.0089,-0.7003) -- (-0.0178,-0.7006) -- (-0.0267,-0.7012) -- (-0.0356,-0.7019) -- (-0.0444,-0.7029) -- (-0.0533,-0.7041) -- (-0.0621,-0.7055) -- (-0.0709,-0.7072) -- (-0.0797,-0.7091) -- (-0.0885,-0.7111) -- (-0.0973,-0.7134) -- (-0.1060,-0.7160) -- (-0.1147,-0.7187) -- (-0.1233,-0.7217) -- (-0.1319,-0.7250) -- (-0.1405,-0.7284) -- (-0.1490,-0.7321) -- (-0.1575,-0.7360) -- (-0.1659,-0.7402) -- (-0.1742,-0.7446) -- (-0.1825,-0.7493) -- (-0.1907,-0.7542) -- (-0.1988,-0.7593) -- (-0.2068,-0.7647) -- (-0.2148,-0.7703) -- (-0.2226,-0.7762) -- (-0.2304,-0.7824) -- (-0.2380,-0.7888) -- (-0.2455,-0.7955) -- (-0.2529,-0.8025) -- (-0.2602,-0.8097) -- (-0.2673,-0.8172) -- (-0.2743,-0.8249) -- (-0.2811,-0.8330) -- (-0.2877,-0.8413) -- (-0.2942,-0.8499) -- (-0.3004,-0.8587) -- (-0.3065,-0.8679) -- (-0.3123,-0.8773) -- (-0.3179,-0.8870) -- (-0.3233,-0.8970) -- (-0.3284,-0.9073) -- (-0.3332,-0.9178) -- (-0.3378,-0.9286) -- (-0.3420,-0.9397);
\draw[wulffB] (0.5000,0.8660) -- (0.4911,0.8512) -- (0.4820,0.8368) -- (0.4726,0.8230) -- (0.4629,0.8097) -- (0.4531,0.7968) -- (0.4430,0.7845) -- (0.4328,0.7726) -- (0.4224,0.7611) -- (0.4119,0.7501) -- (0.4012,0.7396) -- (0.3905,0.7294) -- (0.3796,0.7197) -- (0.3686,0.7103) -- (0.3575,0.7014) -- (0.3464,0.6928) -- (0.3352,0.6846) -- (0.3239,0.6768) -- (0.3126,0.6693) -- (0.3013,0.6622) -- (0.2899,0.6554) -- (0.2784,0.6489) -- (0.2670,0.6428) -- (0.2555,0.6369) -- (0.2439,0.6314) -- (0.2324,0.6262) -- (0.2208,0.6212) -- (0.2092,0.6166) -- (0.1977,0.6122) -- (0.1861,0.6082) -- (0.1745,0.6043) -- (0.1628,0.6008) -- (0.1512,0.5975) -- (0.1396,0.5945) -- (0.1280,0.5917) -- (0.1163,0.5892) -- (0.1047,0.5869) -- (0.0931,0.5849) -- (0.0814,0.5831) -- (0.0698,0.5816) -- (0.0582,0.5803) -- (0.0465,0.5792) -- (0.0349,0.5784) -- (0.0233,0.5778) -- (0.0116,0.5775) -- (0.0000,0.5774) -- (-0.0116,0.5775) -- (-0.0233,0.5778) -- (-0.0349,0.5784) -- (-0.0465,0.5792) -- (-0.0582,0.5803) -- (-0.0698,0.5816) -- (-0.0814,0.5831) -- (-0.0931,0.5849) -- (-0.1047,0.5869) -- (-0.1163,0.5892) -- (-0.1280,0.5917) -- (-0.1396,0.5945) -- (-0.1512,0.5975) -- (-0.1628,0.6008) -- (-0.1745,0.6043) -- (-0.1861,0.6082) -- (-0.1977,0.6122) -- (-0.2092,0.6166) -- (-0.2208,0.6212) -- (-0.2324,0.6262) -- (-0.2439,0.6314) -- (-0.2555,0.6369) -- (-0.2670,0.6428) -- (-0.2784,0.6489) -- (-0.2899,0.6554) -- (-0.3013,0.6622) -- (-0.3126,0.6693) -- (-0.3239,0.6768) -- (-0.3352,0.6846) -- (-0.3464,0.6928) -- (-0.3575,0.7014) -- (-0.3686,0.7103) -- (-0.3796,0.7197) -- (-0.3905,0.7294) -- (-0.4012,0.7396) -- (-0.4119,0.7501) -- (-0.4224,0.7611) -- (-0.4328,0.7726) -- (-0.4430,0.7845) -- (-0.4531,0.7968) -- (-0.4629,0.8097) -- (-0.4726,0.8230) -- (-0.4820,0.8368) -- (-0.4911,0.8512) -- (-0.5000,0.8660);
\draw[wulffB] (0.5000,-0.8660) -- (0.4911,-0.8512) -- (0.4820,-0.8368) -- (0.4726,-0.8230) -- (0.4629,-0.8097) -- (0.4531,-0.7968) -- (0.4430,-0.7845) -- (0.4328,-0.7726) -- (0.4224,-0.7611) -- (0.4119,-0.7501) -- (0.4012,-0.7396) -- (0.3905,-0.7294) -- (0.3796,-0.7197) -- (0.3686,-0.7103) -- (0.3575,-0.7014) -- (0.3464,-0.6928) -- (0.3352,-0.6846) -- (0.3239,-0.6768) -- (0.3126,-0.6693) -- (0.3013,-0.6622) -- (0.2899,-0.6554) -- (0.2784,-0.6489) -- (0.2670,-0.6428) -- (0.2555,-0.6369) -- (0.2439,-0.6314) -- (0.2324,-0.6262) -- (0.2208,-0.6212) -- (0.2092,-0.6166) -- (0.1977,-0.6122) -- (0.1861,-0.6082) -- (0.1745,-0.6043) -- (0.1628,-0.6008) -- (0.1512,-0.5975) -- (0.1396,-0.5945) -- (0.1280,-0.5917) -- (0.1163,-0.5892) -- (0.1047,-0.5869) -- (0.0931,-0.5849) -- (0.0814,-0.5831) -- (0.0698,-0.5816) -- (0.0582,-0.5803) -- (0.0465,-0.5792) -- (0.0349,-0.5784) -- (0.0233,-0.5778) -- (0.0116,-0.5775) -- (0.0000,-0.5774) -- (-0.0116,-0.5775) -- (-0.0233,-0.5778) -- (-0.0349,-0.5784) -- (-0.0465,-0.5792) -- (-0.0582,-0.5803) -- (-0.0698,-0.5816) -- (-0.0814,-0.5831) -- (-0.0931,-0.5849) -- (-0.1047,-0.5869) -- (-0.1163,-0.5892) -- (-0.1280,-0.5917) -- (-0.1396,-0.5945) -- (-0.1512,-0.5975) -- (-0.1628,-0.6008) -- (-0.1745,-0.6043) -- (-0.1861,-0.6082) -- (-0.1977,-0.6122) -- (-0.2092,-0.6166) -- (-0.2208,-0.6212) -- (-0.2324,-0.6262) -- (-0.2439,-0.6314) -- (-0.2555,-0.6369) -- (-0.2670,-0.6428) -- (-0.2784,-0.6489) -- (-0.2899,-0.6554) -- (-0.3013,-0.6622) -- (-0.3126,-0.6693) -- (-0.3239,-0.6768) -- (-0.3352,-0.6846) -- (-0.3464,-0.6928) -- (-0.3575,-0.7014) -- (-0.3686,-0.7103) -- (-0.3796,-0.7197) -- (-0.3905,-0.7294) -- (-0.4012,-0.7396) -- (-0.4119,-0.7501) -- (-0.4224,-0.7611) -- (-0.4328,-0.7726) -- (-0.4430,-0.7845) -- (-0.4531,-0.7968) -- (-0.4629,-0.8097) -- (-0.4726,-0.8230) -- (-0.4820,-0.8368) -- (-0.4911,-0.8512) -- (-0.5000,-0.8660);
\draw[wulffA] (0.6428,0.7660) -- (0.6283,0.7492) -- (0.6137,0.7332) -- (0.5990,0.7178) -- (0.5843,0.7031) -- (0.5695,0.6891) -- (0.5546,0.6757) -- (0.5398,0.6630) -- (0.5249,0.6507) -- (0.5100,0.6391) -- (0.4952,0.6280) -- (0.4803,0.6174) -- (0.4655,0.6073) -- (0.4507,0.5976) -- (0.4360,0.5885) -- (0.4213,0.5797) -- (0.4066,0.5714) -- (0.3920,0.5635) -- (0.3774,0.5560) -- (0.3629,0.5488) -- (0.3484,0.5421) -- (0.3340,0.5357) -- (0.3197,0.5296) -- (0.3053,0.5238) -- (0.2911,0.5184) -- (0.2769,0.5133) -- (0.2627,0.5085) -- (0.2486,0.5040) -- (0.2345,0.4997) -- (0.2205,0.4958) -- (0.2065,0.4921) -- (0.1925,0.4887) -- (0.1786,0.4855) -- (0.1647,0.4826) -- (0.1509,0.4800) -- (0.1371,0.4776) -- (0.1233,0.4754) -- (0.1095,0.4735) -- (0.0958,0.4718) -- (0.0821,0.4703) -- (0.0684,0.4691) -- (0.0547,0.4681) -- (0.0410,0.4673) -- (0.0273,0.4668) -- (0.0137,0.4664) -- (0.0000,0.4663) -- (-0.0137,0.4664) -- (-0.0273,0.4668) -- (-0.0410,0.4673) -- (-0.0547,0.4681) -- (-0.0684,0.4691) -- (-0.0821,0.4703) -- (-0.0958,0.4718) -- (-0.1095,0.4735) -- (-0.1233,0.4754) -- (-0.1371,0.4776) -- (-0.1509,0.4800) -- (-0.1647,0.4826) -- (-0.1786,0.4855) -- (-0.1925,0.4887) -- (-0.2065,0.4921) -- (-0.2205,0.4958) -- (-0.2345,0.4997) -- (-0.2486,0.5040) -- (-0.2627,0.5085) -- (-0.2769,0.5133) -- (-0.2911,0.5184) -- (-0.3053,0.5238) -- (-0.3197,0.5296) -- (-0.3340,0.5357) -- (-0.3484,0.5421) -- (-0.3629,0.5488) -- (-0.3774,0.5560) -- (-0.3920,0.5635) -- (-0.4066,0.5714) -- (-0.4213,0.5797) -- (-0.4360,0.5885) -- (-0.4507,0.5976) -- (-0.4655,0.6073) -- (-0.4803,0.6174) -- (-0.4952,0.6280) -- (-0.5100,0.6391) -- (-0.5249,0.6507) -- (-0.5398,0.6630) -- (-0.5546,0.6757) -- (-0.5695,0.6891) -- (-0.5843,0.7031) -- (-0.5990,0.7178) -- (-0.6137,0.7332) -- (-0.6283,0.7492) -- (-0.6428,0.7660);
\draw[wulffA] (0.6428,-0.7660) -- (0.6283,-0.7492) -- (0.6137,-0.7332) -- (0.5990,-0.7178) -- (0.5843,-0.7031) -- (0.5695,-0.6891) -- (0.5546,-0.6757) -- (0.5398,-0.6630) -- (0.5249,-0.6507) -- (0.5100,-0.6391) -- (0.4952,-0.6280) -- (0.4803,-0.6174) -- (0.4655,-0.6073) -- (0.4507,-0.5976) -- (0.4360,-0.5885) -- (0.4213,-0.5797) -- (0.4066,-0.5714) -- (0.3920,-0.5635) -- (0.3774,-0.5560) -- (0.3629,-0.5488) -- (0.3484,-0.5421) -- (0.3340,-0.5357) -- (0.3197,-0.5296) -- (0.3053,-0.5238) -- (0.2911,-0.5184) -- (0.2769,-0.5133) -- (0.2627,-0.5085) -- (0.2486,-0.5040) -- (0.2345,-0.4997) -- (0.2205,-0.4958) -- (0.2065,-0.4921) -- (0.1925,-0.4887) -- (0.1786,-0.4855) -- (0.1647,-0.4826) -- (0.1509,-0.4800) -- (0.1371,-0.4776) -- (0.1233,-0.4754) -- (0.1095,-0.4735) -- (0.0958,-0.4718) -- (0.0821,-0.4703) -- (0.0684,-0.4691) -- (0.0547,-0.4681) -- (0.0410,-0.4673) -- (0.0273,-0.4668) -- (0.0137,-0.4664) -- (0.0000,-0.4663) -- (-0.0137,-0.4664) -- (-0.0273,-0.4668) -- (-0.0410,-0.4673) -- (-0.0547,-0.4681) -- (-0.0684,-0.4691) -- (-0.0821,-0.4703) -- (-0.0958,-0.4718) -- (-0.1095,-0.4735) -- (-0.1233,-0.4754) -- (-0.1371,-0.4776) -- (-0.1509,-0.4800) -- (-0.1647,-0.4826) -- (-0.1786,-0.4855) -- (-0.1925,-0.4887) -- (-0.2065,-0.4921) -- (-0.2205,-0.4958) -- (-0.2345,-0.4997) -- (-0.2486,-0.5040) -- (-0.2627,-0.5085) -- (-0.2769,-0.5133) -- (-0.2911,-0.5184) -- (-0.3053,-0.5238) -- (-0.3197,-0.5296) -- (-0.3340,-0.5357) -- (-0.3484,-0.5421) -- (-0.3629,-0.5488) -- (-0.3774,-0.5560) -- (-0.3920,-0.5635) -- (-0.4066,-0.5714) -- (-0.4213,-0.5797) -- (-0.4360,-0.5885) -- (-0.4507,-0.5976) -- (-0.4655,-0.6073) -- (-0.4803,-0.6174) -- (-0.4952,-0.6280) -- (-0.5100,-0.6391) -- (-0.5249,-0.6507) -- (-0.5398,-0.6630) -- (-0.5546,-0.6757) -- (-0.5695,-0.6891) -- (-0.5843,-0.7031) -- (-0.5990,-0.7178) -- (-0.6137,-0.7332) -- (-0.6283,-0.7492) -- (-0.6428,-0.7660);
\draw[wulffA] (0.7660,0.6428) -- (0.7456,0.6261) -- (0.7254,0.6102) -- (0.7054,0.5951) -- (0.6855,0.5809) -- (0.6658,0.5673) -- (0.6464,0.5545) -- (0.6271,0.5423) -- (0.6080,0.5307) -- (0.5891,0.5198) -- (0.5704,0.5093) -- (0.5519,0.4995) -- (0.5336,0.4901) -- (0.5154,0.4812) -- (0.4975,0.4728) -- (0.4797,0.4648) -- (0.4621,0.4572) -- (0.4446,0.4500) -- (0.4273,0.4432) -- (0.4102,0.4368) -- (0.3932,0.4307) -- (0.3764,0.4250) -- (0.3597,0.4195) -- (0.3431,0.4144) -- (0.3266,0.4096) -- (0.3103,0.4051) -- (0.2941,0.4008) -- (0.2780,0.3968) -- (0.2620,0.3931) -- (0.2461,0.3897) -- (0.2303,0.3864) -- (0.2145,0.3834) -- (0.1989,0.3807) -- (0.1833,0.3782) -- (0.1678,0.3758) -- (0.1523,0.3737) -- (0.1369,0.3719) -- (0.1216,0.3702) -- (0.1063,0.3687) -- (0.0910,0.3675) -- (0.0758,0.3664) -- (0.0606,0.3655) -- (0.0454,0.3648) -- (0.0303,0.3644) -- (0.0151,0.3641) -- (0.0000,0.3640) -- (-0.0151,0.3641) -- (-0.0303,0.3644) -- (-0.0454,0.3648) -- (-0.0606,0.3655) -- (-0.0758,0.3664) -- (-0.0910,0.3675) -- (-0.1063,0.3687) -- (-0.1216,0.3702) -- (-0.1369,0.3719) -- (-0.1523,0.3737) -- (-0.1678,0.3758) -- (-0.1833,0.3782) -- (-0.1989,0.3807) -- (-0.2145,0.3834) -- (-0.2303,0.3864) -- (-0.2461,0.3897) -- (-0.2620,0.3931) -- (-0.2780,0.3968) -- (-0.2941,0.4008) -- (-0.3103,0.4051) -- (-0.3266,0.4096) -- (-0.3431,0.4144) -- (-0.3597,0.4195) -- (-0.3764,0.4250) -- (-0.3932,0.4307) -- (-0.4102,0.4368) -- (-0.4273,0.4432) -- (-0.4446,0.4500) -- (-0.4621,0.4572) -- (-0.4797,0.4648) -- (-0.4975,0.4728) -- (-0.5154,0.4812) -- (-0.5336,0.4901) -- (-0.5519,0.4995) -- (-0.5704,0.5093) -- (-0.5891,0.5198) -- (-0.6080,0.5307) -- (-0.6271,0.5423) -- (-0.6464,0.5545) -- (-0.6658,0.5673) -- (-0.6855,0.5809) -- (-0.7054,0.5951) -- (-0.7254,0.6102) -- (-0.7456,0.6261) -- (-0.7660,0.6428);
\draw[wulffA] (0.7660,-0.6428) -- (0.7456,-0.6261) -- (0.7254,-0.6102) -- (0.7054,-0.5951) -- (0.6855,-0.5809) -- (0.6658,-0.5673) -- (0.6464,-0.5545) -- (0.6271,-0.5423) -- (0.6080,-0.5307) -- (0.5891,-0.5198) -- (0.5704,-0.5093) -- (0.5519,-0.4995) -- (0.5336,-0.4901) -- (0.5154,-0.4812) -- (0.4975,-0.4728) -- (0.4797,-0.4648) -- (0.4621,-0.4572) -- (0.4446,-0.4500) -- (0.4273,-0.4432) -- (0.4102,-0.4368) -- (0.3932,-0.4307) -- (0.3764,-0.4250) -- (0.3597,-0.4195) -- (0.3431,-0.4144) -- (0.3266,-0.4096) -- (0.3103,-0.4051) -- (0.2941,-0.4008) -- (0.2780,-0.3968) -- (0.2620,-0.3931) -- (0.2461,-0.3897) -- (0.2303,-0.3864) -- (0.2145,-0.3834) -- (0.1989,-0.3807) -- (0.1833,-0.3782) -- (0.1678,-0.3758) -- (0.1523,-0.3737) -- (0.1369,-0.3719) -- (0.1216,-0.3702) -- (0.1063,-0.3687) -- (0.0910,-0.3675) -- (0.0758,-0.3664) -- (0.0606,-0.3655) -- (0.0454,-0.3648) -- (0.0303,-0.3644) -- (0.0151,-0.3641) -- (0.0000,-0.3640) -- (-0.0151,-0.3641) -- (-0.0303,-0.3644) -- (-0.0454,-0.3648) -- (-0.0606,-0.3655) -- (-0.0758,-0.3664) -- (-0.0910,-0.3675) -- (-0.1063,-0.3687) -- (-0.1216,-0.3702) -- (-0.1369,-0.3719) -- (-0.1523,-0.3737) -- (-0.1678,-0.3758) -- (-0.1833,-0.3782) -- (-0.1989,-0.3807) -- (-0.2145,-0.3834) -- (-0.2303,-0.3864) -- (-0.2461,-0.3897) -- (-0.2620,-0.3931) -- (-0.2780,-0.3968) -- (-0.2941,-0.4008) -- (-0.3103,-0.4051) -- (-0.3266,-0.4096) -- (-0.3431,-0.4144) -- (-0.3597,-0.4195) -- (-0.3764,-0.4250) -- (-0.3932,-0.4307) -- (-0.4102,-0.4368) -- (-0.4273,-0.4432) -- (-0.4446,-0.4500) -- (-0.4621,-0.4572) -- (-0.4797,-0.4648) -- (-0.4975,-0.4728) -- (-0.5154,-0.4812) -- (-0.5336,-0.4901) -- (-0.5519,-0.4995) -- (-0.5704,-0.5093) -- (-0.5891,-0.5198) -- (-0.6080,-0.5307) -- (-0.6271,-0.5423) -- (-0.6464,-0.5545) -- (-0.6658,-0.5673) -- (-0.6855,-0.5809) -- (-0.7054,-0.5951) -- (-0.7254,-0.6102) -- (-0.7456,-0.6261) -- (-0.7660,-0.6428);
\draw[wulffB] (0.8660,0.5000) -- (0.8401,0.4853) -- (0.8147,0.4715) -- (0.7898,0.4585) -- (0.7654,0.4462) -- (0.7414,0.4346) -- (0.7178,0.4237) -- (0.6947,0.4134) -- (0.6721,0.4036) -- (0.6498,0.3944) -- (0.6278,0.3857) -- (0.6063,0.3775) -- (0.5851,0.3698) -- (0.5642,0.3624) -- (0.5436,0.3555) -- (0.5234,0.3489) -- (0.5034,0.3427) -- (0.4837,0.3369) -- (0.4643,0.3313) -- (0.4451,0.3261) -- (0.4262,0.3212) -- (0.4075,0.3166) -- (0.3890,0.3122) -- (0.3707,0.3081) -- (0.3526,0.3042) -- (0.3347,0.3006) -- (0.3169,0.2972) -- (0.2993,0.2940) -- (0.2819,0.2910) -- (0.2646,0.2883) -- (0.2474,0.2857) -- (0.2304,0.2833) -- (0.2135,0.2812) -- (0.1967,0.2791) -- (0.1799,0.2773) -- (0.1633,0.2757) -- (0.1467,0.2742) -- (0.1303,0.2729) -- (0.1138,0.2717) -- (0.0975,0.2707) -- (0.0812,0.2699) -- (0.0649,0.2692) -- (0.0486,0.2686) -- (0.0324,0.2683) -- (0.0162,0.2680) -- (0.0000,0.2679) -- (-0.0162,0.2680) -- (-0.0324,0.2683) -- (-0.0486,0.2686) -- (-0.0649,0.2692) -- (-0.0812,0.2699) -- (-0.0975,0.2707) -- (-0.1138,0.2717) -- (-0.1303,0.2729) -- (-0.1467,0.2742) -- (-0.1633,0.2757) -- (-0.1799,0.2773) -- (-0.1967,0.2791) -- (-0.2135,0.2812) -- (-0.2304,0.2833) -- (-0.2474,0.2857) -- (-0.2646,0.2883) -- (-0.2819,0.2910) -- (-0.2993,0.2940) -- (-0.3169,0.2972) -- (-0.3347,0.3006) -- (-0.3526,0.3042) -- (-0.3707,0.3081) -- (-0.3890,0.3122) -- (-0.4075,0.3166) -- (-0.4262,0.3212) -- (-0.4451,0.3261) -- (-0.4643,0.3313) -- (-0.4837,0.3369) -- (-0.5034,0.3427) -- (-0.5234,0.3489) -- (-0.5436,0.3555) -- (-0.5642,0.3624) -- (-0.5851,0.3698) -- (-0.6063,0.3775) -- (-0.6278,0.3857) -- (-0.6498,0.3944) -- (-0.6721,0.4036) -- (-0.6947,0.4134) -- (-0.7178,0.4237) -- (-0.7414,0.4346) -- (-0.7654,0.4462) -- (-0.7898,0.4585) -- (-0.8147,0.4715) -- (-0.8401,0.4853) -- (-0.8660,0.5000);
\draw[wulffB] (0.8660,-0.5000) -- (0.8401,-0.4853) -- (0.8147,-0.4715) -- (0.7898,-0.4585) -- (0.7654,-0.4462) -- (0.7414,-0.4346) -- (0.7178,-0.4237) -- (0.6947,-0.4134) -- (0.6721,-0.4036) -- (0.6498,-0.3944) -- (0.6278,-0.3857) -- (0.6063,-0.3775) -- (0.5851,-0.3698) -- (0.5642,-0.3624) -- (0.5436,-0.3555) -- (0.5234,-0.3489) -- (0.5034,-0.3427) -- (0.4837,-0.3369) -- (0.4643,-0.3313) -- (0.4451,-0.3261) -- (0.4262,-0.3212) -- (0.4075,-0.3166) -- (0.3890,-0.3122) -- (0.3707,-0.3081) -- (0.3526,-0.3042) -- (0.3347,-0.3006) -- (0.3169,-0.2972) -- (0.2993,-0.2940) -- (0.2819,-0.2910) -- (0.2646,-0.2883) -- (0.2474,-0.2857) -- (0.2304,-0.2833) -- (0.2135,-0.2812) -- (0.1967,-0.2791) -- (0.1799,-0.2773) -- (0.1633,-0.2757) -- (0.1467,-0.2742) -- (0.1303,-0.2729) -- (0.1138,-0.2717) -- (0.0975,-0.2707) -- (0.0812,-0.2699) -- (0.0649,-0.2692) -- (0.0486,-0.2686) -- (0.0324,-0.2683) -- (0.0162,-0.2680) -- (0.0000,-0.2679) -- (-0.0162,-0.2680) -- (-0.0324,-0.2683) -- (-0.0486,-0.2686) -- (-0.0649,-0.2692) -- (-0.0812,-0.2699) -- (-0.0975,-0.2707) -- (-0.1138,-0.2717) -- (-0.1303,-0.2729) -- (-0.1467,-0.2742) -- (-0.1633,-0.2757) -- (-0.1799,-0.2773) -- (-0.1967,-0.2791) -- (-0.2135,-0.2812) -- (-0.2304,-0.2833) -- (-0.2474,-0.2857) -- (-0.2646,-0.2883) -- (-0.2819,-0.2910) -- (-0.2993,-0.2940) -- (-0.3169,-0.2972) -- (-0.3347,-0.3006) -- (-0.3526,-0.3042) -- (-0.3707,-0.3081) -- (-0.3890,-0.3122) -- (-0.4075,-0.3166) -- (-0.4262,-0.3212) -- (-0.4451,-0.3261) -- (-0.4643,-0.3313) -- (-0.4837,-0.3369) -- (-0.5034,-0.3427) -- (-0.5234,-0.3489) -- (-0.5436,-0.3555) -- (-0.5642,-0.3624) -- (-0.5851,-0.3698) -- (-0.6063,-0.3775) -- (-0.6278,-0.3857) -- (-0.6498,-0.3944) -- (-0.6721,-0.4036) -- (-0.6947,-0.4134) -- (-0.7178,-0.4237) -- (-0.7414,-0.4346) -- (-0.7654,-0.4462) -- (-0.7898,-0.4585) -- (-0.8147,-0.4715) -- (-0.8401,-0.4853) -- (-0.8660,-0.5000);
\draw[wulffA] (0.9397,0.3420) -- (0.9093,0.3312) -- (0.8797,0.3210) -- (0.8510,0.3114) -- (0.8229,0.3025) -- (0.7956,0.2940) -- (0.7689,0.2861) -- (0.7429,0.2787) -- (0.7175,0.2717) -- (0.6926,0.2651) -- (0.6683,0.2588) -- (0.6444,0.2530) -- (0.6211,0.2474) -- (0.5982,0.2422) -- (0.5757,0.2373) -- (0.5537,0.2327) -- (0.5320,0.2283) -- (0.5107,0.2242) -- (0.4897,0.2203) -- (0.4691,0.2167) -- (0.4488,0.2132) -- (0.4287,0.2100) -- (0.4090,0.2069) -- (0.3895,0.2041) -- (0.3702,0.2014) -- (0.3512,0.1989) -- (0.3324,0.1965) -- (0.3138,0.1943) -- (0.2954,0.1922) -- (0.2771,0.1903) -- (0.2590,0.1886) -- (0.2411,0.1869) -- (0.2233,0.1854) -- (0.2057,0.1840) -- (0.1881,0.1828) -- (0.1707,0.1816) -- (0.1533,0.1806) -- (0.1361,0.1797) -- (0.1189,0.1789) -- (0.1018,0.1782) -- (0.0847,0.1776) -- (0.0677,0.1772) -- (0.0508,0.1768) -- (0.0338,0.1765) -- (0.0169,0.1764) -- (0.0000,0.1763) -- (-0.0169,0.1764) -- (-0.0338,0.1765) -- (-0.0508,0.1768) -- (-0.0677,0.1772) -- (-0.0847,0.1776) -- (-0.1018,0.1782) -- (-0.1189,0.1789) -- (-0.1361,0.1797) -- (-0.1533,0.1806) -- (-0.1707,0.1816) -- (-0.1881,0.1828) -- (-0.2057,0.1840) -- (-0.2233,0.1854) -- (-0.2411,0.1869) -- (-0.2590,0.1886) -- (-0.2771,0.1903) -- (-0.2954,0.1922) -- (-0.3138,0.1943) -- (-0.3324,0.1965) -- (-0.3512,0.1989) -- (-0.3702,0.2014) -- (-0.3895,0.2041) -- (-0.4090,0.2069) -- (-0.4287,0.2100) -- (-0.4488,0.2132) -- (-0.4691,0.2167) -- (-0.4897,0.2203) -- (-0.5107,0.2242) -- (-0.5320,0.2283) -- (-0.5537,0.2327) -- (-0.5757,0.2373) -- (-0.5982,0.2422) -- (-0.6211,0.2474) -- (-0.6444,0.2530) -- (-0.6683,0.2588) -- (-0.6926,0.2651) -- (-0.7175,0.2717) -- (-0.7429,0.2787) -- (-0.7689,0.2861) -- (-0.7956,0.2940) -- (-0.8229,0.3025) -- (-0.8510,0.3114) -- (-0.8797,0.3210) -- (-0.9093,0.3312) -- (-0.9397,0.3420);
\draw[wulffA] (0.9397,-0.3420) -- (0.9093,-0.3312) -- (0.8797,-0.3210) -- (0.8510,-0.3114) -- (0.8229,-0.3025) -- (0.7956,-0.2940) -- (0.7689,-0.2861) -- (0.7429,-0.2787) -- (0.7175,-0.2717) -- (0.6926,-0.2651) -- (0.6683,-0.2588) -- (0.6444,-0.2530) -- (0.6211,-0.2474) -- (0.5982,-0.2422) -- (0.5757,-0.2373) -- (0.5537,-0.2327) -- (0.5320,-0.2283) -- (0.5107,-0.2242) -- (0.4897,-0.2203) -- (0.4691,-0.2167) -- (0.4488,-0.2132) -- (0.4287,-0.2100) -- (0.4090,-0.2069) -- (0.3895,-0.2041) -- (0.3702,-0.2014) -- (0.3512,-0.1989) -- (0.3324,-0.1965) -- (0.3138,-0.1943) -- (0.2954,-0.1922) -- (0.2771,-0.1903) -- (0.2590,-0.1886) -- (0.2411,-0.1869) -- (0.2233,-0.1854) -- (0.2057,-0.1840) -- (0.1881,-0.1828) -- (0.1707,-0.1816) -- (0.1533,-0.1806) -- (0.1361,-0.1797) -- (0.1189,-0.1789) -- (0.1018,-0.1782) -- (0.0847,-0.1776) -- (0.0677,-0.1772) -- (0.0508,-0.1768) -- (0.0338,-0.1765) -- (0.0169,-0.1764) -- (0.0000,-0.1763) -- (-0.0169,-0.1764) -- (-0.0338,-0.1765) -- (-0.0508,-0.1768) -- (-0.0677,-0.1772) -- (-0.0847,-0.1776) -- (-0.1018,-0.1782) -- (-0.1189,-0.1789) -- (-0.1361,-0.1797) -- (-0.1533,-0.1806) -- (-0.1707,-0.1816) -- (-0.1881,-0.1828) -- (-0.2057,-0.1840) -- (-0.2233,-0.1854) -- (-0.2411,-0.1869) -- (-0.2590,-0.1886) -- (-0.2771,-0.1903) -- (-0.2954,-0.1922) -- (-0.3138,-0.1943) -- (-0.3324,-0.1965) -- (-0.3512,-0.1989) -- (-0.3702,-0.2014) -- (-0.3895,-0.2041) -- (-0.4090,-0.2069) -- (-0.4287,-0.2100) -- (-0.4488,-0.2132) -- (-0.4691,-0.2167) -- (-0.4897,-0.2203) -- (-0.5107,-0.2242) -- (-0.5320,-0.2283) -- (-0.5537,-0.2327) -- (-0.5757,-0.2373) -- (-0.5982,-0.2422) -- (-0.6211,-0.2474) -- (-0.6444,-0.2530) -- (-0.6683,-0.2588) -- (-0.6926,-0.2651) -- (-0.7175,-0.2717) -- (-0.7429,-0.2787) -- (-0.7689,-0.2861) -- (-0.7956,-0.2940) -- (-0.8229,-0.3025) -- (-0.8510,-0.3114) -- (-0.8797,-0.3210) -- (-0.9093,-0.3312) -- (-0.9397,-0.3420);
\draw[wulffA] (0.9848,0.1736) -- (0.9515,0.1679) -- (0.9193,0.1625) -- (0.8880,0.1574) -- (0.8577,0.1527) -- (0.8282,0.1483) -- (0.7996,0.1441) -- (0.7717,0.1402) -- (0.7446,0.1366) -- (0.7181,0.1331) -- (0.6923,0.1299) -- (0.6670,0.1269) -- (0.6424,0.1240) -- (0.6182,0.1213) -- (0.5946,0.1187) -- (0.5715,0.1164) -- (0.5488,0.1141) -- (0.5265,0.1120) -- (0.5046,0.1100) -- (0.4831,0.1081) -- (0.4620,0.1063) -- (0.4412,0.1047) -- (0.4206,0.1031) -- (0.4004,0.1016) -- (0.3805,0.1003) -- (0.3608,0.0990) -- (0.3414,0.0978) -- (0.3222,0.0966) -- (0.3032,0.0956) -- (0.2844,0.0946) -- (0.2658,0.0937) -- (0.2473,0.0929) -- (0.2290,0.0921) -- (0.2109,0.0914) -- (0.1928,0.0908) -- (0.1749,0.0902) -- (0.1571,0.0897) -- (0.1394,0.0892) -- (0.1218,0.0888) -- (0.1043,0.0884) -- (0.0868,0.0882) -- (0.0694,0.0879) -- (0.0520,0.0877) -- (0.0347,0.0876) -- (0.0173,0.0875) -- (0.0000,0.0875) -- (-0.0173,0.0875) -- (-0.0347,0.0876) -- (-0.0520,0.0877) -- (-0.0694,0.0879) -- (-0.0868,0.0882) -- (-0.1043,0.0884) -- (-0.1218,0.0888) -- (-0.1394,0.0892) -- (-0.1571,0.0897) -- (-0.1749,0.0902) -- (-0.1928,0.0908) -- (-0.2109,0.0914) -- (-0.2290,0.0921) -- (-0.2473,0.0929) -- (-0.2658,0.0937) -- (-0.2844,0.0946) -- (-0.3032,0.0956) -- (-0.3222,0.0966) -- (-0.3414,0.0978) -- (-0.3608,0.0990) -- (-0.3805,0.1003) -- (-0.4004,0.1016) -- (-0.4206,0.1031) -- (-0.4412,0.1047) -- (-0.4620,0.1063) -- (-0.4831,0.1081) -- (-0.5046,0.1100) -- (-0.5265,0.1120) -- (-0.5488,0.1141) -- (-0.5715,0.1164) -- (-0.5946,0.1187) -- (-0.6182,0.1213) -- (-0.6424,0.1240) -- (-0.6670,0.1269) -- (-0.6923,0.1299) -- (-0.7181,0.1331) -- (-0.7446,0.1366) -- (-0.7717,0.1402) -- (-0.7996,0.1441) -- (-0.8282,0.1483) -- (-0.8577,0.1527) -- (-0.8880,0.1574) -- (-0.9193,0.1625) -- (-0.9515,0.1679) -- (-0.9848,0.1736);
\draw[wulffA] (0.9848,-0.1736) -- (0.9515,-0.1679) -- (0.9193,-0.1625) -- (0.8880,-0.1574) -- (0.8577,-0.1527) -- (0.8282,-0.1483) -- (0.7996,-0.1441) -- (0.7717,-0.1402) -- (0.7446,-0.1366) -- (0.7181,-0.1331) -- (0.6923,-0.1299) -- (0.6670,-0.1269) -- (0.6424,-0.1240) -- (0.6182,-0.1213) -- (0.5946,-0.1187) -- (0.5715,-0.1164) -- (0.5488,-0.1141) -- (0.5265,-0.1120) -- (0.5046,-0.1100) -- (0.4831,-0.1081) -- (0.4620,-0.1063) -- (0.4412,-0.1047) -- (0.4206,-0.1031) -- (0.4004,-0.1016) -- (0.3805,-0.1003) -- (0.3608,-0.0990) -- (0.3414,-0.0978) -- (0.3222,-0.0966) -- (0.3032,-0.0956) -- (0.2844,-0.0946) -- (0.2658,-0.0937) -- (0.2473,-0.0929) -- (0.2290,-0.0921) -- (0.2109,-0.0914) -- (0.1928,-0.0908) -- (0.1749,-0.0902) -- (0.1571,-0.0897) -- (0.1394,-0.0892) -- (0.1218,-0.0888) -- (0.1043,-0.0884) -- (0.0868,-0.0882) -- (0.0694,-0.0879) -- (0.0520,-0.0877) -- (0.0347,-0.0876) -- (0.0173,-0.0875) -- (0.0000,-0.0875) -- (-0.0173,-0.0875) -- (-0.0347,-0.0876) -- (-0.0520,-0.0877) -- (-0.0694,-0.0879) -- (-0.0868,-0.0882) -- (-0.1043,-0.0884) -- (-0.1218,-0.0888) -- (-0.1394,-0.0892) -- (-0.1571,-0.0897) -- (-0.1749,-0.0902) -- (-0.1928,-0.0908) -- (-0.2109,-0.0914) -- (-0.2290,-0.0921) -- (-0.2473,-0.0929) -- (-0.2658,-0.0937) -- (-0.2844,-0.0946) -- (-0.3032,-0.0956) -- (-0.3222,-0.0966) -- (-0.3414,-0.0978) -- (-0.3608,-0.0990) -- (-0.3805,-0.1003) -- (-0.4004,-0.1016) -- (-0.4206,-0.1031) -- (-0.4412,-0.1047) -- (-0.4620,-0.1063) -- (-0.4831,-0.1081) -- (-0.5046,-0.1100) -- (-0.5265,-0.1120) -- (-0.5488,-0.1141) -- (-0.5715,-0.1164) -- (-0.5946,-0.1187) -- (-0.6182,-0.1213) -- (-0.6424,-0.1240) -- (-0.6670,-0.1269) -- (-0.6923,-0.1299) -- (-0.7181,-0.1331) -- (-0.7446,-0.1366) -- (-0.7717,-0.1402) -- (-0.7996,-0.1441) -- (-0.8282,-0.1483) -- (-0.8577,-0.1527) -- (-0.8880,-0.1574) -- (-0.9193,-0.1625) -- (-0.9515,-0.1679) -- (-0.9848,-0.1736);
\draw[wulffB] (-1,0) -- (1,0) (0,-1) -- (0,1);
\draw[wulffA] (1.0000,0.0000) -- (0.9400,0.0000);
\draw[wulffA] (0.9994,0.0349) -- (0.9644,0.0337);
\draw[wulffA] (0.9976,0.0698) -- (0.9626,0.0673);
\draw[wulffA] (0.9945,0.1045) -- (0.9597,0.1009);
\draw[wulffA] (0.9903,0.1392) -- (0.9556,0.1343);
\draw[wulffA] (0.9848,0.1736) -- (0.9257,0.1632);
\draw[wulffA] (0.9781,0.2079) -- (0.9439,0.2006);
\draw[wulffA] (0.9703,0.2419) -- (0.9363,0.2335);
\draw[wulffA] (0.9613,0.2756) -- (0.9276,0.2660);
\draw[wulffA] (0.9511,0.3090) -- (0.9178,0.2982);
\draw[wulffA] (0.9397,0.3420) -- (0.8833,0.3215);
\draw[wulffA] (0.9272,0.3746) -- (0.8947,0.3615);
\draw[wulffA] (0.9135,0.4067) -- (0.8816,0.3925);
\draw[wulffA] (0.8988,0.4384) -- (0.8673,0.4230);
\draw[wulffA] (0.8829,0.4695) -- (0.8520,0.4530);
\draw[wulffA] (0.8660,0.5000) -- (0.8141,0.4700);
\draw[wulffA] (0.8480,0.5299) -- (0.8184,0.5114);
\draw[wulffA] (0.8290,0.5592) -- (0.8000,0.5396);
\draw[wulffA] (0.8090,0.5878) -- (0.7807,0.5672);
\draw[wulffA] (0.7880,0.6157) -- (0.7604,0.5941);
\draw[wulffA] (0.7660,0.6428) -- (0.7201,0.6042);
\draw[wulffA] (0.7431,0.6691) -- (0.7171,0.6457);
\draw[wulffA] (0.7193,0.6947) -- (0.6942,0.6703);
\draw[wulffA] (0.6947,0.7193) -- (0.6703,0.6942);
\draw[wulffA] (0.6691,0.7431) -- (0.6457,0.7171);
\draw[wulffA] (0.6428,0.7660) -- (0.6042,0.7201);
\draw[wulffA] (0.6157,0.7880) -- (0.5941,0.7604);
\draw[wulffA] (0.5878,0.8090) -- (0.5672,0.7807);
\draw[wulffA] (0.5592,0.8290) -- (0.5396,0.8000);
\draw[wulffA] (0.5299,0.8480) -- (0.5114,0.8184);
\draw[wulffA] (0.5000,0.8660) -- (0.4700,0.8141);
\draw[wulffA] (0.4695,0.8829) -- (0.4530,0.8520);
\draw[wulffA] (0.4384,0.8988) -- (0.4230,0.8673);
\draw[wulffA] (0.4067,0.9135) -- (0.3925,0.8816);
\draw[wulffA] (0.3746,0.9272) -- (0.3615,0.8947);
\draw[wulffA] (0.3420,0.9397) -- (0.3215,0.8833);
\draw[wulffA] (0.3090,0.9511) -- (0.2982,0.9178);
\draw[wulffA] (0.2756,0.9613) -- (0.2660,0.9276);
\draw[wulffA] (0.2419,0.9703) -- (0.2335,0.9363);
\draw[wulffA] (0.2079,0.9781) -- (0.2006,0.9439);
\draw[wulffA] (0.1736,0.9848) -- (0.1632,0.9257);
\draw[wulffA] (0.1392,0.9903) -- (0.1343,0.9556);
\draw[wulffA] (0.1045,0.9945) -- (0.1009,0.9597);
\draw[wulffA] (0.0698,0.9976) -- (0.0673,0.9626);
\draw[wulffA] (0.0349,0.9994) -- (0.0337,0.9644);
\draw[wulffA] (0.0000,1.0000) -- (0.0000,0.9400);
\draw[wulffA] (-0.0349,0.9994) -- (-0.0337,0.9644);
\draw[wulffA] (-0.0698,0.9976) -- (-0.0673,0.9626);
\draw[wulffA] (-0.1045,0.9945) -- (-0.1009,0.9597);
\draw[wulffA] (-0.1392,0.9903) -- (-0.1343,0.9556);
\draw[wulffA] (-0.1736,0.9848) -- (-0.1632,0.9257);
\draw[wulffA] (-0.2079,0.9781) -- (-0.2006,0.9439);
\draw[wulffA] (-0.2419,0.9703) -- (-0.2335,0.9363);
\draw[wulffA] (-0.2756,0.9613) -- (-0.2660,0.9276);
\draw[wulffA] (-0.3090,0.9511) -- (-0.2982,0.9178);
\draw[wulffA] (-0.3420,0.9397) -- (-0.3215,0.8833);
\draw[wulffA] (-0.3746,0.9272) -- (-0.3615,0.8947);
\draw[wulffA] (-0.4067,0.9135) -- (-0.3925,0.8816);
\draw[wulffA] (-0.4384,0.8988) -- (-0.4230,0.8673);
\draw[wulffA] (-0.4695,0.8829) -- (-0.4530,0.8520);
\draw[wulffA] (-0.5000,0.8660) -- (-0.4700,0.8141);
\draw[wulffA] (-0.5299,0.8480) -- (-0.5114,0.8184);
\draw[wulffA] (-0.5592,0.8290) -- (-0.5396,0.8000);
\draw[wulffA] (-0.5878,0.8090) -- (-0.5672,0.7807);
\draw[wulffA] (-0.6157,0.7880) -- (-0.5941,0.7604);
\draw[wulffA] (-0.6428,0.7660) -- (-0.6042,0.7201);
\draw[wulffA] (-0.6691,0.7431) -- (-0.6457,0.7171);
\draw[wulffA] (-0.6947,0.7193) -- (-0.6703,0.6942);
\draw[wulffA] (-0.7193,0.6947) -- (-0.6942,0.6703);
\draw[wulffA] (-0.7431,0.6691) -- (-0.7171,0.6457);
\draw[wulffA] (-0.7660,0.6428) -- (-0.7201,0.6042);
\draw[wulffA] (-0.7880,0.6157) -- (-0.7604,0.5941);
\draw[wulffA] (-0.8090,0.5878) -- (-0.7807,0.5672);
\draw[wulffA] (-0.8290,0.5592) -- (-0.8000,0.5396);
\draw[wulffA] (-0.8480,0.5299) -- (-0.8184,0.5114);
\draw[wulffA] (-0.8660,0.5000) -- (-0.8141,0.4700);
\draw[wulffA] (-0.8829,0.4695) -- (-0.8520,0.4530);
\draw[wulffA] (-0.8988,0.4384) -- (-0.8673,0.4230);
\draw[wulffA] (-0.9135,0.4067) -- (-0.8816,0.3925);
\draw[wulffA] (-0.9272,0.3746) -- (-0.8947,0.3615);
\draw[wulffA] (-0.9397,0.3420) -- (-0.8833,0.3215);
\draw[wulffA] (-0.9511,0.3090) -- (-0.9178,0.2982);
\draw[wulffA] (-0.9613,0.2756) -- (-0.9276,0.2660);
\draw[wulffA] (-0.9703,0.2419) -- (-0.9363,0.2335);
\draw[wulffA] (-0.9781,0.2079) -- (-0.9439,0.2006);
\draw[wulffA] (-0.9848,0.1736) -- (-0.9257,0.1632);
\draw[wulffA] (-0.9903,0.1392) -- (-0.9556,0.1343);
\draw[wulffA] (-0.9945,0.1045) -- (-0.9597,0.1009);
\draw[wulffA] (-0.9976,0.0698) -- (-0.9626,0.0673);
\draw[wulffA] (-0.9994,0.0349) -- (-0.9644,0.0337);
\draw[wulffA] (-1.0000,0.0000) -- (-0.9400,0.0000);
\draw[wulffA] (-0.9994,-0.0349) -- (-0.9644,-0.0337);
\draw[wulffA] (-0.9976,-0.0698) -- (-0.9626,-0.0673);
\draw[wulffA] (-0.9945,-0.1045) -- (-0.9597,-0.1009);
\draw[wulffA] (-0.9903,-0.1392) -- (-0.9556,-0.1343);
\draw[wulffA] (-0.9848,-0.1736) -- (-0.9257,-0.1632);
\draw[wulffA] (-0.9781,-0.2079) -- (-0.9439,-0.2006);
\draw[wulffA] (-0.9703,-0.2419) -- (-0.9363,-0.2335);
\draw[wulffA] (-0.9613,-0.2756) -- (-0.9276,-0.2660);
\draw[wulffA] (-0.9511,-0.3090) -- (-0.9178,-0.2982);
\draw[wulffA] (-0.9397,-0.3420) -- (-0.8833,-0.3215);
\draw[wulffA] (-0.9272,-0.3746) -- (-0.8947,-0.3615);
\draw[wulffA] (-0.9135,-0.4067) -- (-0.8816,-0.3925);
\draw[wulffA] (-0.8988,-0.4384) -- (-0.8673,-0.4230);
\draw[wulffA] (-0.8829,-0.4695) -- (-0.8520,-0.4530);
\draw[wulffA] (-0.8660,-0.5000) -- (-0.8141,-0.4700);
\draw[wulffA] (-0.8480,-0.5299) -- (-0.8184,-0.5114);
\draw[wulffA] (-0.8290,-0.5592) -- (-0.8000,-0.5396);
\draw[wulffA] (-0.8090,-0.5878) -- (-0.7807,-0.5672);
\draw[wulffA] (-0.7880,-0.6157) -- (-0.7604,-0.5941);
\draw[wulffA] (-0.7660,-0.6428) -- (-0.7201,-0.6042);
\draw[wulffA] (-0.7431,-0.6691) -- (-0.7171,-0.6457);
\draw[wulffA] (-0.7193,-0.6947) -- (-0.6942,-0.6703);
\draw[wulffA] (-0.6947,-0.7193) -- (-0.6703,-0.6942);
\draw[wulffA] (-0.6691,-0.7431) -- (-0.6457,-0.7171);
\draw[wulffA] (-0.6428,-0.7660) -- (-0.6042,-0.7201);
\draw[wulffA] (-0.6157,-0.7880) -- (-0.5941,-0.7604);
\draw[wulffA] (-0.5878,-0.8090) -- (-0.5672,-0.7807);
\draw[wulffA] (-0.5592,-0.8290) -- (-0.5396,-0.8000);
\draw[wulffA] (-0.5299,-0.8480) -- (-0.5114,-0.8184);
\draw[wulffA] (-0.5000,-0.8660) -- (-0.4700,-0.8141);
\draw[wulffA] (-0.4695,-0.8829) -- (-0.4530,-0.8520);
\draw[wulffA] (-0.4384,-0.8988) -- (-0.4230,-0.8673);
\draw[wulffA] (-0.4067,-0.9135) -- (-0.3925,-0.8816);
\draw[wulffA] (-0.3746,-0.9272) -- (-0.3615,-0.8947);
\draw[wulffA] (-0.3420,-0.9397) -- (-0.3215,-0.8833);
\draw[wulffA] (-0.3090,-0.9511) -- (-0.2982,-0.9178);
\draw[wulffA] (-0.2756,-0.9613) -- (-0.2660,-0.9276);
\draw[wulffA] (-0.2419,-0.9703) -- (-0.2335,-0.9363);
\draw[wulffA] (-0.2079,-0.9781) -- (-0.2006,-0.9439);
\draw[wulffA] (-0.1736,-0.9848) -- (-0.1632,-0.9257);
\draw[wulffA] (-0.1392,-0.9903) -- (-0.1343,-0.9556);
\draw[wulffA] (-0.1045,-0.9945) -- (-0.1009,-0.9597);
\draw[wulffA] (-0.0698,-0.9976) -- (-0.0673,-0.9626);
\draw[wulffA] (-0.0349,-0.9994) -- (-0.0337,-0.9644);
\draw[wulffA] (-0.0000,-1.0000) -- (-0.0000,-0.9400);
\draw[wulffA] (0.0349,-0.9994) -- (0.0337,-0.9644);
\draw[wulffA] (0.0698,-0.9976) -- (0.0673,-0.9626);
\draw[wulffA] (0.1045,-0.9945) -- (0.1009,-0.9597);
\draw[wulffA] (0.1392,-0.9903) -- (0.1343,-0.9556);
\draw[wulffA] (0.1736,-0.9848) -- (0.1632,-0.9257);
\draw[wulffA] (0.2079,-0.9781) -- (0.2006,-0.9439);
\draw[wulffA] (0.2419,-0.9703) -- (0.2335,-0.9363);
\draw[wulffA] (0.2756,-0.9613) -- (0.2660,-0.9276);
\draw[wulffA] (0.3090,-0.9511) -- (0.2982,-0.9178);
\draw[wulffA] (0.3420,-0.9397) -- (0.3215,-0.8833);
\draw[wulffA] (0.3746,-0.9272) -- (0.3615,-0.8947);
\draw[wulffA] (0.4067,-0.9135) -- (0.3925,-0.8816);
\draw[wulffA] (0.4384,-0.8988) -- (0.4230,-0.8673);
\draw[wulffA] (0.4695,-0.8829) -- (0.4530,-0.8520);
\draw[wulffA] (0.5000,-0.8660) -- (0.4700,-0.8141);
\draw[wulffA] (0.5299,-0.8480) -- (0.5114,-0.8184);
\draw[wulffA] (0.5592,-0.8290) -- (0.5396,-0.8000);
\draw[wulffA] (0.5878,-0.8090) -- (0.5672,-0.7807);
\draw[wulffA] (0.6157,-0.7880) -- (0.5941,-0.7604);
\draw[wulffA] (0.6428,-0.7660) -- (0.6042,-0.7201);
\draw[wulffA] (0.6691,-0.7431) -- (0.6457,-0.7171);
\draw[wulffA] (0.6947,-0.7193) -- (0.6703,-0.6942);
\draw[wulffA] (0.7193,-0.6947) -- (0.6942,-0.6703);
\draw[wulffA] (0.7431,-0.6691) -- (0.7171,-0.6457);
\draw[wulffA] (0.7660,-0.6428) -- (0.7201,-0.6042);
\draw[wulffA] (0.7880,-0.6157) -- (0.7604,-0.5941);
\draw[wulffA] (0.8090,-0.5878) -- (0.7807,-0.5672);
\draw[wulffA] (0.8290,-0.5592) -- (0.8000,-0.5396);
\draw[wulffA] (0.8480,-0.5299) -- (0.8184,-0.5114);
\draw[wulffA] (0.8660,-0.5000) -- (0.8141,-0.4700);
\draw[wulffA] (0.8829,-0.4695) -- (0.8520,-0.4530);
\draw[wulffA] (0.8988,-0.4384) -- (0.8673,-0.4230);
\draw[wulffA] (0.9135,-0.4067) -- (0.8816,-0.3925);
\draw[wulffA] (0.9272,-0.3746) -- (0.8947,-0.3615);
\draw[wulffA] (0.9397,-0.3420) -- (0.8833,-0.3215);
\draw[wulffA] (0.9511,-0.3090) -- (0.9178,-0.2982);
\draw[wulffA] (0.9613,-0.2756) -- (0.9276,-0.2660);
\draw[wulffA] (0.9703,-0.2419) -- (0.9363,-0.2335);
\draw[wulffA] (0.9781,-0.2079) -- (0.9439,-0.2006);
\draw[wulffA] (0.9848,-0.1736) -- (0.9257,-0.1632);
\draw[wulffA] (0.9903,-0.1392) -- (0.9556,-0.1343);
\draw[wulffA] (0.9945,-0.1045) -- (0.9597,-0.1009);
\draw[wulffA] (0.9976,-0.0698) -- (0.9626,-0.0673);
\draw[wulffA] (0.9994,-0.0349) -- (0.9644,-0.0337);
\draw[wulffC] (0,0) circle (1);

  \node[font=\scriptsize, anchor=south] at (0,1.02) {$y$};
\end{tikzpicture}
\caption{}
\end{subfigure}\hfill
\begin{subfigure}[b]{0.48\textwidth}\centering
\begin{tikzpicture}[scale=2.9]
  % создано tools/gen_stereo.py: сетка Болдырева (r = 1)
\draw[wulffA] (0,0) circle (0.0875);
\draw[wulffA] (0,0) circle (0.1763);
\draw[wulffB] (0,0) circle (0.2679);
\draw[wulffA] (0,0) circle (0.3640);
\draw[wulffA] (0,0) circle (0.4663);
\draw[wulffB] (0,0) circle (0.5774);
\draw[wulffA] (0,0) circle (0.7002);
\draw[wulffA] (0,0) circle (0.8391);
\draw[wulffB] (1.0000,0.0000) -- (-1.0000,-0.0000);
\draw[wulffA] (0.9848,0.1736) -- (-0.9848,-0.1736);
\draw[wulffA] (0.9397,0.3420) -- (-0.9397,-0.3420);
\draw[wulffB] (0.8660,0.5000) -- (-0.8660,-0.5000);
\draw[wulffA] (0.7660,0.6428) -- (-0.7660,-0.6428);
\draw[wulffA] (0.6428,0.7660) -- (-0.6428,-0.7660);
\draw[wulffB] (0.5000,0.8660) -- (-0.5000,-0.8660);
\draw[wulffA] (0.3420,0.9397) -- (-0.3420,-0.9397);
\draw[wulffA] (0.1736,0.9848) -- (-0.1736,-0.9848);
\draw[wulffB] (0.0000,1.0000) -- (-0.0000,-1.0000);
\draw[wulffA] (-0.1736,0.9848) -- (0.1736,-0.9848);
\draw[wulffA] (-0.3420,0.9397) -- (0.3420,-0.9397);
\draw[wulffB] (-0.5000,0.8660) -- (0.5000,-0.8660);
\draw[wulffA] (-0.6428,0.7660) -- (0.6428,-0.7660);
\draw[wulffA] (-0.7660,0.6428) -- (0.7660,-0.6428);
\draw[wulffB] (-0.8660,0.5000) -- (0.8660,-0.5000);
\draw[wulffA] (-0.9397,0.3420) -- (0.9397,-0.3420);
\draw[wulffA] (-0.9848,0.1736) -- (0.9848,-0.1736);
\draw[wulffC] (0,0) circle (1);

\end{tikzpicture}
\caption{}
\end{subfigure}
\caption{Координатные сетки (расчёт): \textit{а}~— сетка Вульфа (меридианы
и параллели через $10^\circ$, жирные~— через $30^\circ$; штрихи на основном
круге~— через $2^\circ$); \textit{б}~— сетка Болдырева (окружности равного
угла $\theta$ с осью $z$ и радиусы через $10^\circ$)}
\label{fig:nets}
\end{figure}

\begin{podrobnee}[Как построить меридиан и параллель]
Меридиан, изображающий плоскость, наклонённую под углом $\delta$ к
плоскости проекции (и проходящую через ось $y$),~— дуга окружности
радиуса $R_\text{м}=r/\cos\delta$ с центром на экваторе на расстоянии
$OO_\text{м}=r\operatorname{tg}\delta$ от центра. Параллель~— малый круг угловым
радиусом $\beta$ вокруг оси $y$~— дуга окружности радиуса
$R_\text{п}=r\operatorname{tg}\beta$ с центром на оси $y$ на расстоянии
$OO_\text{п}=r/\cos\beta$ (рис.~\ref{fig:merid}). Обе формулы следуют из
\eqref{eq:stereo-rho} и того, что окружность проходит через известные
точки основного круга\isprmark.
\end{podrobnee}
\isprtext{В оригинале (формулы (1.4)--(1.7)) меридиан характеризуется углом
$\alpha=90^\circ-\delta$, а параллель~— угловым расстоянием от экватора
$90^\circ-\beta$; в этих обозначениях $R_\text{м}=r/\sin\alpha$,
$OO_\text{м}=r\operatorname{ctg}\alpha$, $R_\text{п}=r\operatorname{ctg}\beta'$,
$OO_\text{п}=r/\sin\beta'$. Здесь использованы углы, отсчитываемые от
плоскости проекции и от оси сетки,~— так формулы проще запомнить.}

\begin{figure}[ht]
\centering
\begin{tikzpicture}[scale=1.0]
  \def\R{2.2}
  \draw[atomLine, thick] (0,0) circle (\R);
  \draw[fmGray] (-3.2,0) -- (2.5,0) (0,-2.5) -- (0,3.6);
  \fill (0,0) circle (1.2pt) node[below right, font=\small]{$O$};
  \begin{scope}
    \clip (0,0) circle (\R);
    \draw[fmBlue, very thick] (-2.622,0) circle (3.4227);
    \draw[fmRed, very thick] (0,2.872) circle (1.846);
  \end{scope}
  \fill[fmBlue] (-2.622,0) circle (1.4pt) node[below, font=\small]{$O_\text{м}$};
  \draw[fmBlue, dashed] (-2.622,0) -- (0,2.2) node[midway, above left=-2pt, font=\small]{$R_\text{м}$};
  \fill[fmRed] (0,2.872) circle (1.4pt) node[right, font=\small]{$O_\text{п}$};
  \draw[fmRed, dashed] (0,2.872) -- ({2.2*sin(40)},{2.2*cos(40)}) node[pos=0.45, right, font=\small]{$R_\text{п}$};
  \fill (0,2.2) circle (1.2pt) node[above left, font=\small]{$y$};
  \node[fmBlue, font=\small] at (0.75,-1.3) {меридиан};
  \node[fmRed, font=\small] at (0,0.75) {параллель};
\end{tikzpicture}
\caption{Построение меридиана ($\delta=50^\circ$, синий) и параллели
($\beta=40^\circ$, красный): центры $O_\text{м}$ и $O_\text{п}$ лежат вне основного
круга}
\label{fig:merid}
\end{figure}

На практике проекцию строят на кальке, наложенной на сетку того же
диаметра: кальку поворачивают вокруг общего центра (следя, чтобы центр не
смещался) до нужного положения и считывают углы по сетке.

\paragraph{Основные задачи, решаемые с помощью сеток}
\begin{enumerate}
\item \emph{Угол между двумя направлениями.} Два направления комплекса
всегда лежат в одной плоскости, а плоскость на проекции~— меридиан.
Поэтому кальку поворачивают, пока обе точки не окажутся на одном
меридиане, и отсчитывают угол между ними вдоль меридиана (по
пересекаемым параллелям). Если одно направление изображено точкой, а
другое крестиком, их выводят на меридианы, симметричные относительно
нулевого, и складывают углы, отсчитанные до полюса сетки. Можно и
заменить крестик противоположной ему точкой (отложив равный отрезок по
другую сторону от центра) и найти дополнительный угол $180^\circ-\alpha$.
\item \emph{Координаты проекции по углам.} Если направление образует с
осью $z$ угол $\theta$, а его проекция на плоскость $xy$~— угол $\varphi$
с осью $x$, то
\begin{equation}\label{eq:stereo-xy}
R_x=r\operatorname{tg}\frac\theta2\cos\varphi,\qquad R_y=r\operatorname{tg}\frac\theta2\sin\varphi .
\end{equation}
Через направляющие углы $\alpha$, $\beta$, $\gamma$ с осями $x$, $y$, $z$:
$R_x=r\operatorname{tg}(\gamma/2)\cos\alpha/\sin\gamma$,
$R_y=r\operatorname{tg}(\gamma/2)\cos\beta/\sin\gamma$.
\item \emph{Полюс плоскости.} Плоскость~— меридиан; поворачиваем кальку до
совпадения её дуги с меридианом сетки и отсчитываем по экватору
$90^\circ$: полученная точка~— проекция нормали, т.\,е. полюс плоскости.
\item \emph{Угол между плоскостями}~— это угол между их нормалями:
находим полюсы и измеряем угол между ними, как в задаче~1.
\item \emph{Поворот вокруг оси, лежащей в плоскости проекции.} Совмещаем
ось поворота с полюсом сетки; при повороте каждая точка движется по своей
\emph{параллели} (по конусу вокруг оси) на заданный угол. Точка, «перешедшая»
через основной круг, уходит в южное полушарие и изображается крестиком
(её можно перевести обратно в точку отражением через центр). Плоскость
поворачивают через её полюс: поворачивают полюс, затем строят новую дугу.
\item \emph{Поворот вокруг оси, перпендикулярной плоскости проекции}~—
по сетке Болдырева: все точки смещаются по концентрическим окружностям на
один и тот же угол.
\item \emph{Малый круг}~— геометрическое место направлений, образующих угол
$\alpha$ с заданной осью (проекция кругового конуса). Через проекцию оси
проводят меридиан, откладывают по нему в обе стороны $\alpha$ и через две
полученные точки проводят окружность. Её центр \emph{не совпадает} с
проекцией оси из-за неравномерности шкалы. Поворот вокруг \emph{произвольной}
оси переводит точку по такому малому кругу; угол поворота откладывают не
транспортиром, а в плоскости, перпендикулярной оси.
\end{enumerate}

\begin{primer}[Проекция направления {[111]} кубического кристалла]
Плоскость проекции~— $(001)$. Направление $[111]$ образует с $[001]$ угол
$\theta=\arccos(1/\sqrt3)=54{,}74^\circ$, а его проекция на плоскость $xy$
идёт под $\varphi=45^\circ$ к оси $x$. По~\eqref{eq:stereo-xy}
$\rho=r\operatorname{tg}27{,}37^\circ=0{,}518\,r$, $R_x=R_y=0{,}366\,r$. Угол между
$[111]$ и $[001]$, измеренный по сетке Вульфа (обе точки выведены на
нулевой меридиан), равен тем же $54{,}7^\circ$.
\end{primer}

\paragraph{Описание поворота: углы Эйлера и ось--угол}
Поворот твёрдого тела (кристалла) задают двумя способами.
\begin{itemize}
\item \emph{Углы Эйлера} $\varphi$, $\theta$, $\psi$: поворот на $\varphi$ вокруг
оси $z$ (оси $x,y$ переходят в $x',y'$), затем на $\theta$ вокруг нового
положения оси $x'$ (ось $z$ переходит в $z'$), затем на $\psi$ вокруг $z'$.
По проекции с осями до и после поворота углы находят так: $\theta$~— угол
между $z$ и $z'$; $\varphi$~— угол, на котором линия пересечения плоскостей
$xy$ и $x'y'$ (ось $x'$ промежуточной системы) отстоит от $x$; $\psi$~— угол
от этой линии до конечной оси $x''$.
\item \emph{Ось и угол поворота} (чаще используются в кристаллографии,
например для описания разориентации зёрен, п.~1.7). Если направления $1$
и $2$ переходят в $1'$ и $2'$, ось поворота лежит на пересечении двух
\emph{биссекторных} плоскостей: каждая проходит через биссектрису угла
между направлением до и после поворота и через нормаль к плоскости,
содержащей эти два направления. Угол поворота~— центральный угол в
плоскости, перпендикулярной оси.
\end{itemize}
Матричная запись поворотов~— в п.~\ref{ssec:orth}.
